\documentclass[11pt]{article}
\usepackage[hmargin=1in,vmargin=.95in]{geometry}
\usepackage{amsmath,amssymb,amsthm,mathtools,booktabs,array,longtable,microtype,xcolor}
\usepackage[colorlinks=true,allcolors=blue!50!black]{hyperref}
\hypersetup{pdfauthor={GPT-6 Astra},pdftitle={Interactads: A Common Calculus for Context and Effects},pdfsubject={A common calculus for context, content and their interaction},pdfkeywords={contextads; contentads; interactads; result-dependent continuations; resource aggregation; relative double categories; decorated displays; bispans; pseudodistributive laws}}
\numberwithin{equation}{section}
\newtheorem{theorem}{Theorem}[section]
\newtheorem{proposition}[theorem]{Proposition}
\newtheorem{lemma}[theorem]{Lemma}
\newtheorem{corollary}[theorem]{Corollary}
\theoremstyle{definition}\newtheorem{definition}[theorem]{Definition}
\newcommand{\B}{\mathcal B}\newcommand{\J}{\mathcal J}
\newcommand{\D}{\mathbb D}\newcommand{\N}{\mathbb N}
\newcommand{\K}{\mathcal K}\newcommand{\Lcat}{\mathcal L}
\newcommand{\V}{\mathcal V}\newcommand{\Set}{\mathsf{Set}}\newcommand{\Cat}{\mathsf{Cat}}
\newcommand{\id}{\mathrm{id}}\newcommand{\op}{\mathrm{op}}\newcommand{\dom}{\mathrm{dom}}\newcommand{\cod}{\mathrm{cod}}
\newcommand{\iso}{\xrightarrow{\sim}}\newcommand{\ea}{e}\newcommand{\un}{\mathbf1}\newcommand{\Tt}{\overline T}
\newcommand{\Disc}{\mathrm{disc}}\newcommand{\st}{\mathbin{\star}}\newcommand{\di}{\mathbin{\diamond}}\newcommand{\bx}{\mathop{\boxtimes}}
\newcommand{\pure}{\mathsf p}\newcommand{\grade}{\mathsf G}
\title{Interactads: A Common Calculus for Context and Effects}
\author{GPT-6 Astra}\date{21 November 2026}
\newcommand{\Prob}{\mathbf{Prob}^{+}_{\mathrm{fin}}}
\newcommand{\List}{\mathsf{List}}
\begin{document}\maketitle
\begin{abstract}
A computation can produce a result that determines what the next stage
requires. Returning two requests instead of one can demand two resource
choices; changing their labels can change the resources from which those
choices are drawn. Composing such computations requires an account of how
produced effects and required context act on one another.

An \emph{interactad} gives a common calculus for this interaction. Building
on Capucci and Myers' contextads, their output-side dual contentads, and
their proposal to relate the two by distributive laws, it combines separately
specified context and content presentations over a shared category of pure
maps. A full pseudodistributive law makes the mixed computations compose
associatively up to coherent isomorphism. The resulting calculus retains
maps of resources and produced content, including noninvertible maps, and
recovers the original one-sided calculi when the other side is specialized
to the identity. Embedding both sides faithfully in a fixed mixed calculus
requires further conditions.

A sufficient construction obtains the interaction from a shared description
of output occurrences, under explicit cartesian lifting, pasting-stability
and pure-interface compatibility hypotheses. In a concrete finite probability/list
construction, each returned position requests a separate draw from the resource
distribution selected by its label. The resulting law lets successive stages
compose coherently even when intermediate results change the numbers and
types of those draws. It also tracks how resource aggregation and maps of
produced content propagate through composition, retaining distinctions that
an output distribution alone can erase.
\end{abstract}
\noindent{\small
Project sources and verification:
Appendix~\ref{app:reuse}. The principal general results have written proofs;
selected subresults, not the full construction, are formalized in Lean.}

\section{Computation that changes its continuation}\label{sec:question}
A composite computation should be usable as a stage in a larger
computation. This is a substantive requirement when one stage's output
determines what its successor needs. The composition rule must account
for those needs in a form that the next composition can still use.

Consider a survey planner that returns a list of sites to sample. The
sampler needs a source of readings for each position in the list. Two
entries for the same site request two readings under independent
sampling; replacing one of the entries changes the source for that
position. The readings may in turn prompt a list of follow-up samples.
Whether the planner is first composed with the sampler, or the sampler
with the follow-up rule, the composite should describe the same sampling
process, down to the correspondence between requests and readings.
Section~\ref{sec:prob} works out this pattern for finite probability spaces.

A request is part of what one stage produces; a source of readings is
part of what another stage requires. That distinction matters also when
comparing computations. Two ways of consolidating result records can
agree on which records merge and still combine their underlying evidence
differently. To compose such comparisons, a calculus needs both the map
on records and the map on resources, with a rule that relates them when
a record specifies new requirements.

Capucci and Myers' contextads unify context-dependent arrows, including
comonadic, parameterized and span constructions; their contentads give
the output-side dual \cite{CM}. Their Section~7 proposes distributive laws
for contextful-contentful maps and shows how some produced effects can be
represented as context. The question here is when separately presented
context and content calculi admit a law for their interaction that also
accounts for their maps and changes of interface.

\paragraph{The unification.}
The common algebra has a pure category $C$ and a category $H(A,B)$ of
computations between each pair of interfaces. It composes relative to the
pure maps, rather than forgetting which maps are pure. A context presentation
has bodies $e_AP\to B$; a content presentation has bodies $A\to d_BQ$.
An interactad combines these presentations with a coherent crossing and has
mixed bodies
\[
                 e_AP\longrightarrow d_BQ.
\]
Here $P$ is the context grade, $Q$ the content grade, and $e,d$ their
extensions. Grades are objects of categories, not merely indices: their maps
can merge witnesses, aggregate coefficients or compare receipts. The common
calculus retains those maps, including noninvertible ones.

The mixed composition rule explains how separately specified input and output
extensions act on one another, and supplies a common place to compare their
specializations and semantics. This goes beyond a one-sided presentation
without ruling out one-sided encodings of particular instances. The required
compatibility of the two sides is part of the result.

\paragraph{Results and reading route.}
Theorem~\ref{thm:unification} assembles the relative construction, the mixed
normal form and its recovered interface squares. Section~\ref{sec:special}
identifies the familiar one-sided specializations. Theorem~\ref{thm:problaw}
constructs a nonthin probability/list interaction and proves that both
one-sided calculi embed locally fully faithfully in that example.
Section~\ref{sec:independent-main} then explains how the general
independent-input theorem constructs, rather than assumes, an interaction:
its complete hypotheses and proof are in Appendices~\ref{app:construction}--\ref{app:parameters}.
Section~\ref{sec:paired} exhibits genuinely varying content grades and
resource maps, so the unification is not restricted to terminal-grade effects.
Section~\ref{sec:observations} separates the native calculus from its
observable comparisons. Section~\ref{sec:limits} establishes a limitation
on conservative assembly and locates further extensions. Readers interested
in the computational mechanism can begin with Section~\ref{sec:prob}; readers
reusing the categorical construction can begin with
Section~\ref{sec:relative} and Appendix~\ref{app:construction}.

\section{A common algebra with a visible pure interface}\label{sec:relative}
Work with a small ordinary category $C$, in nested universes when necessary.
In particular, a category of sets at one universe level is used as a small
base only at a larger level. A category-valued module is a pseudofunctor
$X:C^{\op}\times C\to\Cat$. For two modules the relative tensor is
\begin{equation}\label{eq:relative-tensor}
 (X\odot_CY)(A,D)=\int^{B\in C}_{\mathrm{ps}}X(A,B)\times Y(B,D).
\end{equation}
Tensors are written in execution order. The unit is
$I_C(A,B)=\Disc C(A,B)$. This is the ordinary-base part of Lawler's
category-valued profunctor tricategory, with the variance convention made
explicit in the companion construction \cite[Sec.~4.1.1]{Lawler}\cite{U}.
No new ambient tricategory is proposed here.

An object of the pseudocoend has a representative $[B;x,y]$. Its morphisms
include the native maps in $X(A,B)$ and $Y(B,D)$, and invertible balances
\[
 [B';u_*x,y']\cong[B;x,u^*y']\qquad(u:B\to B').
\]
Balance means that a pure map at an internal interface can be performed on
either side of that interface. The balances satisfy naturality, identity and
composition relations. Their invertibility does not invert the native maps
in the two computation categories. Replacing this tensor by a set of
isomorphism classes would erase precisely the comparisons the paper retains.

\begin{definition}[Relative computation algebra]\label{def:relative-main}
A relative category-valued pseudopromonad is a module $H$, a unit
$j:I_C\to H$, and a multiplication $m:H\odot_CH\to H$, with invertible
associativity and unit modifications satisfying the pseudomonad axioms.
Its concrete realization $B_H$ has hom categories $H(A,B)$ and composition
$m([h,k])$. The designated pure pseudofunctor is $J(u)=j(u)$.
\end{definition}
The same data can be recovered from a bicategory $B$ with an
identity-on-objects pure pseudofunctor $J:C\to B$: set
$H(A,B)=B(A,B)$ and let pure maps act by whiskering with $J$.
The bicategory associator supplies pure-middle balance. This equivalence of
descriptions retains each hom category unchanged; it is not merely the
assertion that several examples are bicategories.

Changes of endpoints remain visible. A square from $h:A\rightsquigarrow B$
to $h':A'\rightsquigarrow B'$ over $u:A\to A'$ and $v:B\to B'$ is
\begin{equation}\label{eq:square-main}
                   J(v)h\Longrightarrow h'J(u).
\end{equation}
These cells form the relative quintet double category. If a pseudo-double
category already has chosen companions for its vertical maps, bending a
square along their binding cells gives exactly \eqref{eq:square-main}.
Bending back cancels the binding cells by the companion equations, so this
is a bijection on squares with each fixed boundary, with pasting preserved
up to the canonical constraints \cite[Sec.~3]{U}. No conjoint is needed,
and noninvertible squares are not lost.

\subsection{The two presentations side by side}\label{sec:presentations-main}
Let $K:C^{\op}\to\Cat$ be a context indexing pseudofunctor with extension
$e:\int K\to C$. A grade $P\in K_A$ extends $A$ to $e_AP$; reindexing
along $u:A'\to A$ has comparison $e_u(P):e_{A'}u^*P\to e_AP$.
Its contextad data include a unit $I_A$ with discharge
$\varepsilon_A:e_AI_A\to A$ and dependent accumulation $P\st R$
for $R\in K_{e_AP}$, with
$\delta_{P,R}:e_A(P\st R)\to e_{e_AP}R$.
The structure maps $\varepsilon,\delta$ need not be invertible. The
horizontal associativity and unit comparisons are invertible and obey the
contextad laws.

Dually, a content indexing pseudofunctor $L:C\to\Cat$ has extension
$d:\int L\to C$, transport $v_!Q$ and comparison
$d_v(Q):d_BQ\to d_{B'}v_!Q$. Its unit $U_B$ has body
$\eta^L_B:B\to d_BU_B$. For $Q\in L_B$ and $S\in L_{d_BQ}$,
accumulation $Q\di S$ has body
$\mu^L_{Q,S}:d_{d_BQ}S\to d_B(Q\di S)$. These bodies need not be
invertible either. The one-sided hom categories are
\begin{equation}\label{eq:onesided-main}
 M(A,B)=(e_A\downarrow B),\qquad N(A,B)=(A\downarrow d_B).
\end{equation}
A context cell $\alpha:(P,f)\to(P',f')$ satisfies $f=f'e_A\alpha$;
a content cell $\beta:(Q,g)\to(Q',g')$ satisfies $d_B\beta\,g=g'$.
Their endpoint actions are
\[
 (P,f)\mapsto(u^*P,vf e_u(P)),\qquad
 (Q,g)\mapsto(v_!Q,d_v(Q)gu).
\]
These formulas fix the directions of all grade maps.

For successive context arrows $(P,f)$ and $(R,g)$, composition restricts
$R$ along $f$, accumulates $P\st f^*R$, and uses the body
\[
 e_A(P\st f^*R)\xrightarrow{\delta}e_{e_AP}f^*R
                  \xrightarrow{e_f(R)}e_BR\xrightarrow g D.
\]
For successive content arrows $(Q,f)$ and $(S,g)$, it transports $Q$
along $g$, accumulates $S\di g_!Q$, and uses
\[
 A\xrightarrow f d_BQ\xrightarrow{d_g(Q)}d_{d_DS}g_!Q
                  \xrightarrow{\mu^L}d_D(S\di g_!Q).
\]
The contextad and contentad laws give relative pseudopromonads $M,N$.
Neither of these formulas yet says how an output extension should change
the next arrow's input extension.

\section{Mixed normal form and its composition}\label{sec:mixed-main}
The missing operation moves content followed by context into context
followed by content. A \emph{full pseudodistributive law}
\begin{equation}\label{eq:law-main}
                 \tau:N\odot_CM\longrightarrow M\odot_CN
\end{equation}
is the appropriate compatibility datum. Its crossing need not be
invertible. Its four comparisons with units and multiplications are
invertible and satisfy the complete coherence axioms of
\cite[Def.~4.1]{GL}, not just their four boundaries. This is established
pseudomonad-composition machinery \cite{marmolejo,street}.

\begin{definition}[Interactad]\label{def:interactad}
An interactad consists of valid context and content presentations over a
specified pure category, their induced modules \eqref{eq:onesided-main},
and a full law \eqref{eq:law-main}. The presentations and the pure
interface are retained as part of the package. An \emph{induced} interactad
also records a derivation of that law from independently stated hypotheses.
\end{definition}
This usage follows the project terminology note linked in Appendix~\ref{app:reuse}. A general
relative pseudopromonad can describe other computational structures without
being presented as this particular two-sided package.

The following result assembles the normalization and mixed-algebra results
of the unified construction \cite{U}.
\begin{theorem}[The common presented calculus]\label{thm:unification}
For context and content presentations as in
Section~\ref{sec:presentations-main}, there is a pseudonatural equivalence
\begin{equation}\label{eq:normal-main}
 M\odot_CN\simeq H,\qquad H(A,B)=(e_A\downarrow d_B).
\end{equation}
It retains every native cell and all endpoint transports. A full law
\eqref{eq:law-main} makes $H$ a relative pseudopromonad with unit and
multiplication transported from
\begin{equation}\label{eq:mult-main}
 I_C\simeq I_C\odot_CI_C\xrightarrow{j_M\odot j_N}MN,
 \qquad MNMN\xrightarrow{M\tau N}MMNN\xrightarrow{m_Mm_N}MN.
\end{equation}
Its realization has mixed arrows $(P,Q,f:e_AP\to d_BQ)$ and native cells
$(\alpha,\beta)$ satisfying $d_B\beta\,f=f'e_A\alpha$. Its endpoint
squares over $(u,v)$ are exactly
\begin{equation}\label{eq:mixedsquare-main}
 \alpha:P\to u^*P',\quad\beta:v_!Q\to Q',\qquad
 d_{B'}\beta\,d_v(Q)f=f'e_u(P')e_A\alpha.
\end{equation}
Specializing the other side to the identity recovers each original
one-sided calculus. Full faithfulness of its unit-induced comparison inside
a fixed mixed calculus is an additional property, not a general consequence.
\end{theorem}
\begin{proof}
At fixed endpoints, send
$[Z;P,x:eP\to Z;Q,y:Z\to dQ]$ to $(P,Q,yx)$, its product cells to
their native pair, and its pure balances to identities. Normalize a mixed
arrow to $[eP;P,\id;Q,f]$. For a mixed cell $(\alpha,\beta)$, its
reverse image is the path
\[
 [eP;P,\id;Q,f]\xrightarrow{\beta}
 [eP;P,\id;Q',d\beta f]
 \xrightarrow{b_{e\alpha}}
 [eP';P,e\alpha;Q',f']\xrightarrow{\alpha}
 [eP';P',\id;Q',f'].
\]
Balance naturality and $b_{vu}=b_vb_u$ show that these paths preserve
composition. The forward composite is the identity. The reverse composite
is isomorphic to the identity by the balance $b_x$ along the factorization
map $x$. Its naturality on product maps and balance generators gives the
required equivalence of categories, with no inversion of $\alpha$ or
$\beta$.

Under endpoint transport the composed body is
$d_v(Q)yx e_u(P)$. The comparison on the reverse normalization is the
specific balance along $e_u(P)$, together with the content transport
comparison. The balance relations and cleavage equations prove
pseudonaturality. This gives \eqref{eq:normal-main} globally, not just a
bijection on arrows. The composite-pseudomonad theorem then gives
\eqref{eq:mult-main}; its associativity and unit modifications transport
along that equivalence. The squares are the native mixed cell equation
after output and input transport, hence \eqref{eq:mixedsquare-main}.
Pseudoco-Yoneda gives $M\odot I_C\simeq M$ and $I_C\odot N\simeq N$.
\end{proof}

The theorem is operational. On generic content and context arrows,
normalize $\tau$ to grades $\widehat R\in K_{d_BQ}$ and
$\widehat Q\in L_{e_BR}$, with
$\chi:e_{d_BQ}\widehat R\to d_{e_BR}\widehat Q$.
For $h=(P,Q,f):A\rightsquigarrow B$ followed by
$k=(R,S,g):B\rightsquigarrow D$, the composite grades are
\begin{equation}\label{eq:grades-main}
                 P\st f^*\widehat R,\qquad S\di g_!\widehat Q.
\end{equation}
Its body is the typed composite
\begin{align*}
 e_A(P\st f^*\widehat R)&\xrightarrow{\delta}e_{e_AP}f^*\widehat R
 \xrightarrow{e_f}e_{d_BQ}\widehat R\xrightarrow\chi d_{e_BR}\widehat Q\\
 &\xrightarrow{d_g}d_{d_DS}g_!\widehat Q
 \xrightarrow{\mu^L}d_D(S\di g_!\widehat Q).
\end{align*}
Thus the first result selects the lifted requirements of the next stage,
while the next body transports the first stage's retained content. The two
sides do not merely occupy adjacent fields in a record.

The pure arrow $u$ has representative $(I_A,U_B,\eta^L_Bu\varepsilon_A)$.
The one-sided maps into this fixed mixed calculus are
\[
 (P,f)\mapsto(P,U_B,\eta^L_Bf),\qquad
 (Q,g)\mapsto(I_A,Q,g\varepsilon_A).
\]
A nonmonic $\eta^L_B$ can identify different context bodies. For example,
the identity context calculus and the terminal monad on sets have the
canonical unit crossing, yet the context comparison sends every function
$A\to B$ to the unique function $A\to 1$. A lawful chosen crossing need
not embed both original calculi conservatively.

\subsection{What presentation recovery requires}
The common operational algebra does not intrinsically identify its supplied
grade categories. Mark the generic context arrows
$\kappa_P=(P,\id):A\rightsquigarrow e_AP$, their squares, factorizations
and chosen cartesian lifts. A square between two such arrows has
$\alpha:P\to u^*P'$ and forces its other boundary to be
$e_u(P')e_A\alpha$. It therefore recovers exactly a morphism of $\int K$.
Marked unit and accumulation factorizations retain $\varepsilon$ and
$\delta$. The dual generic content arrows
$\gamma_Q=(Q,\id):d_BQ\rightsquigarrow B$ recover $\int L$ and its
opfibration in the same way \cite[Sec.~4]{U}. This is recovery of a
supplied marked presentation. Marking only its images after a nonfaithful
mixed comparison cannot restore information already lost.

\section{Familiar computations, recovered at their actual scope}\label{sec:special}
The value of one construction is that its specializations recover the
operations and comparisons of the familiar calculi, rather than just giving
their arrows the same schematic appearance. The following cases are
one-sided specializations of Theorem~\ref{thm:unification}; the identity
calculus supplies the unused side.

\paragraph{Kleisli and coKleisli.}
For a monad $T$, take terminal content grade fibres and $d_B(*)=TB$.
Then $N(A,B)=\Disc C(A,TB)$, with composite $\mu_D T(g)f$ and pure map
$\eta_Bu$. The endpoint-square equation is $T(v)f=f'u$.
For a comonad $G$, terminal context fibres with $e_A(*)=GA$ give
$M(A,B)=\Disc C(GA,B)$, with composite $gG(f)\delta_A$, pure map
$u\varepsilon_A$, and square equation $vf=f'G(u)$.
Terminal grade fibres do not mean trivial effects: the monad or comonad
structure is in the extension and its structure maps.
Given a separate mixed law $\lambda:GT\to TG$, the fixed-effect mixed
specialization has arrows $GA\to TB$ and composite
\[
             \mu_D T(g)\lambda_B G(f)\delta_A.
\]
This is a conditional specialization, not an automatic law between an
arbitrary monad and comonad.

\paragraph{Forward parameterization.}
For a right monoidal action of $V$ on $C$, take constant context fibres
$K_A=V$ and $e_AP=A\mathbin{\odot}P$. Then arrows are
$(P,f:A\odot P\to B)$, and forward reparameterization is
$\alpha:P\to P'$ with $f=f'(1\odot\alpha)$. Composition uses parameter
$P\otimes Q$ and the action associator followed by $f\odot1$ and $g$;
squares satisfy $vf=f'(u\odot\alpha)$. This is the forward-cell convention
for Para. It does not supply copying of arbitrary monoidal parameters,
reverse differentiation or an optimization algorithm. The contextad
unification of this case belongs to Capucci--Myers \cite{CM}.

\paragraph{Display spans and dual cospans.}
Let $\mathcal D$ contain identities and be closed under composition and
pullback, with the requisite chosen pullbacks. Taking
$K_A=\mathcal D/A$ and $e$ to be domain gives spans
$A\xleftarrow p X\xrightarrow f B$ with $p\in\mathcal D$.
Cells are all commuting apex maps, composition uses $X\times_BY$, and
endpoint squares are maps $a:X\to X'$ with $p'a=up$, $f'a=vf$.
The dual construction uses a pushout-stable class of codisplays, chosen
pushouts, and content extension by codomain; it gives the corresponding
cospans and their apex maps. These are exact native recoveries. Their
coexistence in the same ambient construction supplies no general crossing
between display spans and codisplay cospans.

\paragraph{Changing presentation is a theorem, not a slogan.}
Some one-sided forms can be transposed. For $GA=E\times A$, currying
identifies coKleisli arrows with $A\to B^E$ when the exponentials exist;
the endpoint equations and reuse of the same environment agree. Writer
arrows can likewise be presented by discrete context grades recording a
function $w:A\to W$, for a monoid $W$ \cite[Sec.~8]{U}. Such explicit
comparisons explain why context versus content is not a universal
expressivity dichotomy. They do not identify arbitrary nonthin
parameterization with reader or writer computation, and they do not remove
the compatibility problem for independently given presentations.

\section{Dependent probability and lists do work together}\label{sec:prob}
Let $\Prob$ have finite nonempty probability spaces with strictly positive
weights, and deterministic measure-preserving maps. Thus $\alpha:P\to P'$
means $p'(y)=\sum_{\alpha x=y}p(x)$. It retains both parallel maps and
noninvertible folds. Over sets, let
\[
 K_A={\Prob}^A,\qquad e_AP=\coprod_{a\in A}P_a.
\]
Accumulating $P$ and a family $R$ over $e_AP$ gives the fibre
$\coprod_{x\in P_a}R_{a,x}$ with weight $p_a(x)r_{a,x}(y)$.
The unit is a singleton family. For content use terminal grade fibres with
extension $d_B(*)=\List B$, singleton insertion and concatenation.
The content grade category is terminal, but the list content is nontrivial.

A mixed computation is $(P,f:\coprod_aP_a\to\List B)$. Its native cell
$\alpha:(P,f)\to(P',f')$ is measure-preserving and satisfies
$f'(a,\alpha_ax)=f(a,x)$. Under endpoint maps $(u,v)$ the equation becomes
\begin{equation}\label{eq:prob-square-main}
                    f'(u(a),\alpha_ax)=\List(v)f(a,x).
\end{equation}
For a list $\ell=[b_1,\ldots,b_n]$, crossing a continuation family $R$
back through this content gives
\begin{equation}\label{eq:prob-cross-main}
 \widehat R_\ell=\prod_{i=1}^nR_{b_i},\qquad
 \widehat r_\ell(r_1,\ldots,r_n)=\prod_i r_{b_i}(r_i),\qquad
 \chi(\ell,(r_i))=[(b_i,r_i)]_i.
\end{equation}
Every occurrence has its own coordinate, even when labels repeat. The
empty product is a singleton. This is the independent probability tensor,
not a claim that carrier products are categorical products in $\Prob$.

\begin{theorem}[A conservative nonthin mixed instance]\label{thm:problaw}
Formula~\eqref{eq:prob-cross-main} defines a full pseudodistributive law.
For $h=(P,f):A\rightsquigarrow B$ and $k=(R,g):B\rightsquigarrow D$,
write $f(a,x)=[b_1,\ldots,b_n]$. Their composite has
\begin{align}
 \Omega_{k\circ h}(a)&=\coprod_{x\in P_a}\prod_{i=1}^nR_{b_i},
       \label{eq:prob-comp-main}\\
 w(x,(r_i))&=p_a(x)\prod_i r_{b_i}(r_i),\notag\\
 (k\circ h)(a,x,(r_i))&=\operatorname{concat}_i g(b_i,r_i).\notag
\end{align}
It retains all native grade maps and endpoint squares. The unit-induced
context and content comparisons are locally fully faithful.
\end{theorem}
\begin{proof}
The fibres are finite and nonempty, every weight is positive, and their
total is
\[\sum_xp_a(x)\prod_i\sum_{r_i}r_{b_i}(r_i)=1.\]
For maps $\beta_b:R_b\to R'_b$, the product map is measure-preserving
because
\[
 \sum_{\beta_{b_i}r_i=r'_i\ \forall i}\prod_i r_{b_i}(r_i)
   =\prod_i\sum_{\beta_{b_i}r_i=r'_i}r_{b_i}(r_i)
   =\prod_i r'_{b_i}(r'_i).
\]
A pure map relabels positions without merging them, so the crossing respects
balance even for a noninjective pure map. Mixed cells compose by
$(x,(r_i))\mapsto(\alpha_ax,(\beta_{b_i}r_i))$; summing over product
preimages and then over the $\alpha$-preimage gives the target weight.
Equation~\eqref{eq:prob-square-main} proves endpoint compatibility.

For content multiplication the comparison groups a product over a
concatenated list into the products over its ordered blocks. For context
multiplication it is the addressed distributivity bijection
\[
 \prod_i\coprod_{y_i\in R_i}S_{i,y_i}
       \cong\coprod_{(y_i)\in\prod_iR_i}\prod_iS_{i,y_i}.
\]
Both carriers assign the same weight to the same coordinates. Units insert
or delete singleton factors. Label every coordinate by its sequence of list
positions and context stages. Each structural comparison preserves that
address, its value and its weight factor. Every route through a full
law-coherence diagram is therefore the same function; empty blocks delete
only forced singleton choices. This verifies the law, including naturality
on noninvertible grade maps.

For context arrows, the inclusion uses singleton output lists, so its cell
equation is precisely the original context equation. For content arrows,
the inclusion uses singleton sample families, between which there is a
unique map; the mixed equation holds exactly when the list functions agree.
Both inclusions are consequently full and faithful on each hom category,
with the original composition and square comparisons.
\end{proof}

\subsection{A result-dependent calculation}
Take a singleton input, a prefix sample with two equally weighted values
$x_0,x_1$, and outputs $f(x_0)=[b]$, $f(x_1)=[c,c]$. At $b$ the
continuation draws a fair bit and returns its singleton list. At $c$ it
draws a value in $\{0,1,2\}$ with probabilities $(1/2,1/3,1/6)$ and also
returns a singleton list. The mixed context has eleven witnesses: two over
$x_0$, each of weight $1/4$, and nine over $x_1$, whose weights, indexed by
the two trit values, are
\[
 \frac12\begin{pmatrix}1/2\\1/3\\1/6\end{pmatrix}
          \begin{pmatrix}1/2&1/3&1/6\end{pmatrix}.
\]
The first branch returns one bit; the second returns an ordered pair of
trits. The prefix has changed both the number of draws and their probability
fibre. The repeated $c$ labels still give separate occurrence coordinates.
A fixed continuation fibre or a fixed count of copies does not give this
retained witness presentation without further encoding and a comparison.
Both the output shape and the input-dependent continuation family contribute
to \eqref{eq:prob-comp-main}.

For three computations the associator is the explicit address regrouping
\[
 ((x,(y_i)),(z_{ij})_{\rm flattened})
       \longleftrightarrow(x,((y_i,(z_{ij})_j))_i).
\]
The second output determines each block length. The comparison preserves
all samples and their weight factors, including when the lengths differ
between witnesses. This is why the two-sided composition is coherent.

At the level of output laws, the independent draws in
\eqref{eq:prob-cross-main} are the ordinary list/distribution law of
Fujii, Tsai, Montacute and Hasuo \cite[Ex.~20, Prop.~36, Ex.~38 and Thm.~42]{fujii}.
For example, two requests of the same label, each returning $[]$ or $[z]$
with equal probability, yield four addressed outcomes of weight $1/4$.
Concatenation sends two distinct outcomes to $[z]$, producing masses
$1/4,1/2,1/4$ on $[],[z],[z,z]$. Taking the output law folds the two
middle witnesses together by a noninvertible measure-preserving map.
Finite indexed valuations already distinguish samples with the same
observed output, up to weight-preserving renaming
\cite[Defs.~4.1--4.3]{varaccawinskel}. The native calculus also retains
deterministic measure-preserving maps between presentations and
transports them through result-dependent list substitution.

The example also identifies information that a distribution-only account
forgets. A constant output on an equally weighted two-point sample space
has two distinct automorphisms and a noninvertible map to the singleton
presentation. All remain cells of $H$. Their output distributions agree,
but their native maps do not. A distribution semantics may be useful; it
is a subsequent comparison, not the definition of the combined calculus.

\section{Constructing a crossing from separate inputs}\label{sec:independent-main}
The crossing $\tau$ moves a content computation followed by a context
computation into context-then-content order. Here $N$ denotes the content
calculus, $M$ the context calculus, and $\tau:N\odot M\to M\odot N$ must
respect changes at their shared pure interface. The probability/list law
constructs such a crossing by working directly with lists and probability
spaces. The following construction applies whenever both sides can use the
same occurrence geometry.

The base category $\B$ has pullbacks and a class $\J$ of displays containing
isomorphisms and closed under composition, pullback and isomorphism of
arrows. A cartesian monad
$T$ preserves pullbacks and has pullback naturality squares for its unit
and multiplication; it also preserves $\J$. The unit supplies the trivial
continuation and multiplication substitutes one occurrence shape into
another. A context is a display $r:E\to B$ with resource decoration, and
its chosen codomain-cartesian restrictions are stable under horizontal
composition. The \emph{pseudo-double lift} $\Tt$ applies $T$ to the
underlying displays while acting on their resource decorations and maps.
It preserves codomain-cartesian restriction squares and has cartesian unit and
multiplication squares and coherent horizontal comparisons. Content
independently assigns an occurrence map $a_n:A\to TB$ to each arrow
$n:A\rightsquigarrow B$, compatibly with native squares, pure maps and
composition.

Both sides meet on the same object. Lifting $n$ along $r$ must agree on its
carrier with applying the occurrence action to the decorated context and
restricting along $a_n$:
\begin{equation}\label{eq:common-carrier-main}
                     A_r\cong A\times_{TB}TE.
\end{equation}
For lists, $a_n(a)=[b_1,\ldots,b_k]$ means that a point of this pullback over
$a$ is a list $[e_1,\ldots,e_k]$ with $r(e_i)=b_i$ for each $i$: one
allowed resource choice per addressed position, even when labels repeat.
The content square must be genuinely codomain-cartesian, testing all
squares, with horizontal pastes stable along the allowed displays.
Canonical pure pullback squares must also be cartesian. That last
condition supplies the invertible Beck--Chevalley comparison needed for
balance over arbitrary pure maps. It cannot be replaced by the existence
of an isomorphism between carriers.

Under exactly the input contracts of Definitions~\ref{def:lift} and
\ref{def:corners}, Theorem~\ref{thm:glue} constructs the full law with
normalized value
\[
             \tau[n,(P,g:E\to C)]=(a_n^*\Tt P,\ g_*n_r).
\]
The context map is obtained by the lifted occurrence action; the content
cell is obtained by cartesian factorization. They share a carrier map but
remain separately identifiable. The four structural comparisons are
constructed from these same data, and all ten pastings, including their
pure seams, are checked in Appendix~\ref{app:coherence}. In comma
presentations the result returns exactly the native mixed category of
Theorem~\ref{thm:unification}, not a quotient or an alternate presentation
with untracked grade cells (Corollary~\ref{cor:set}).

Finite polynomial occurrences provide a broad resource instance. For a
coherent symmetric rig category $\V$, each finite resource fibre has
coefficient objects $V_e$ and a total augmentation
$\bigoplus_eV_e\to I$. A merge into $e'$ carries an actual map
$\bigoplus_{he=e'}V_e\to V'_{e'}$. At an occurrence shape with positions
$i$, resource choices have rows $\bigotimes_iV_{e_i}$; maps distribute
finite tensors over the preimage sums and apply the original coefficient
maps. The operations are sums and tensors of the chosen rig, not assumed
categorical products and coproducts. Appendix~\ref{app:construction}
constructs this lift and verifies its coherence on directed maps.

Section~\ref{sec:paired} uses this construction to compose varying content
grades and noninvertible receipt maps with resource aggregation.
Appendix~\ref{app:credit} relates the construction to its polynomial,
relative-corner and wreath precursors.

\section{Resource aggregation and content receipts}\label{sec:paired}
A second mixed computation uses a contentad whose grades are
\emph{receipts}. Over a set $B$, a receipt grade is a set $Q$ with a label
$q:Q\to\mathsf{List}(B)$. Its extended result is $dQ=B+Q$: an output
can be an ordinary $B$-value or a receipt describing a list of requested
continuations. On the ordinary summand the occurrence list is the singleton;
on the receipt summand it is $q$. A map of grades preserves these lists
under the specified change of $B$. Transport changes labels, not the set
of receipts. Accumulation keeps both sets of receipts; a label mentioning an
intermediate receipt is flattened by substituting that receipt's own list. Section~\ref{sec:receipts}
checks the full content structure.

Take all three interfaces to be the singleton $1=\{*\}$. An arrow $F$
has trivial context and one receipt $q$ labelled $[*,*]$; its body returns
$q$. An arrow $G$ has two context choices $0,1$, each with vector-space
coefficient object $k=\mathbb Q$. Its augmentation is the linear map
$k\oplus k\to k$, $(a,b)\mapsto a+b$. Its content has one
receipt $s$ labelled $[*]$, and its body returns $s$ on either choice.
In particular the second arrow has nontrivial content too.

The composite first chooses $(x,z)\in\{0,1\}^2$, with coefficient
$k\otimes k$ at that choice. Lifting the receipt $q$ produces four
receipts $q_{xz}$ labelled $[x,z]$. The body of $G$ relabels each of them
by $[s,s]$ without merging them. Accumulation with $s$ then substitutes
its singleton label, so all four have label $[*,*]$. The resulting grade
has five receipts,
\[
                 \{s,q_{00},q_{01},q_{10},q_{11}\}.
\]
The composite body takes $(x,z)$ to $q_{xz}$. The receipt $s$ remains in
the grade although the body does not reach it. Replacing the grade by the
image of the body would change the construction.

For comparison, let $G'$ have one choice with coefficient object $V=k^2$,
augmented
by $\epsilon(a,b)=a+b$, and the same receipt $s$ and constant body. The
fold of choices admits the two coefficient maps
\begin{equation}\label{eq:rankmaps}
 \gamma=1_V,\qquad \gamma'=p,\qquad p(a,b)=(a+b,0).
\end{equation}
They give two native maps $G\to G'$. After composition with $F$, their
coefficient maps, under the addressed identification
$\bigoplus_{x,z}k\cong V\otimes V$, are
\[
                 1_{V\otimes V},\qquad p\otimes p.
\]
The vector $e_2\otimes e_2$ goes respectively to $e_2\otimes e_2$ and
$e_1\otimes e_1$. Their grade map is the same: all four $q_{xz}$ go to
the single new receipt $q'$, while $s$ goes to $s$. This is a genuinely
noninvertible grade map. It is not the relabelling transport, which kept
four tags, but the comparison into the target lifted grade.

Thus there are two distinct pieces of map-level information in one
composite. The coefficient map distinguishes the two aggregations; the
grade map records the four-to-one receipt comparison. Neither can replace
the other. The categorical product of two copies of $(V,\epsilon:V\to k)$ in vector
spaces over $k$ with augmentation to $k$ is the equal-total pullback
\[
 V\times_kV=\{(v,w)\in V\oplus V:\epsilon(v)=\epsilon(w)\}.
\]
It has dimension $3$, since equality of totals imposes one independent
linear equation on $V\oplus V$. The tensor used in the composite has
dimension $4$.


The readout of the second-pair coefficient is an external inspection of
the native coefficient maps, not a claim that all native closing contexts
separate them. A closing map to the augmented unit sees only the total
augmentation.

\section{From native computations to observations}\label{sec:observations}
The common calculus gives a place for comparisons between implementations,
but does not prescribe which observations an application should retain.
The finite probability/list construction has a natural output-law map
\[
 E(P,f)(a)=\sum_{x\in P_a}p_a(x)\,\delta_{f(a,x)}.
\]
In its target, composition samples an outer list and then independently
samples a continuation list at each addressed occurrence before
concatenating. Finite distributivity groups the native witnesses by their
output lists and proves that $E$ preserves composition. This argument
requires the continuation to use its actual interface value, not to inspect
a witness that the map $E$ has erased.

Why retain the law of the whole list? If a result requests $[b,b]$ and each
position returns a fair bit, independent draws yield $00,01,10,11$ with
probability $1/4$ each. Reusing one draw yields $00,11$ with probability
$1/2$ each. The one-position marginals agree, while the probabilities of
disagreement are $1/2$ and $0$. This compares observation models; it does
not provide two laws for the same interaction calculus. The finite
probability/list law uses independent draws. A shared record must be
retained in the relevant input or private state and given its own semantics.

The applications use a different semantic model, still within the same
finite occurrence construction. Finite causal trees replace lists so that
runs retain their event structure, and witnesses carry cost as well as
probability. The corresponding comparison sends a native command to the
joint law of its cost and causal tree:
\begin{equation}\label{eq:joint-main}
 L(t)(a)=\sum_{h\in\Omega_t(a)}p_t(h)
              \delta_{(c_t(h),\theta_t(h))}.
\end{equation}
Trees have finitely many addressed holes; composition grafts the selected
child trees and adds their costs. For a prefix witness $h$ and children
$k_i$ at its actual leaves, mass and cost are
$p_t(h)\prod_i p_u(k_i)$ and $c_t(h)+\sum_i c_u(k_i)$.
Terminal content grades again carry a nontrivial tree monad.
In the stated finite language, the separate evaluator agrees with
the semantics on addressed samples, weights, costs and typed causal trees.
Its run correspondence~\mbox{[\href{https://github.com/eternal-november/eos-research-programmes/blob/10e771644cfdf58c3d05e989e4b26cae7164c80a/projects/contextual-composition/protocols/plural-semantics/Plural_Semantics_for_Result_Dependent_Protocols.pdf}{PDF}, \href{https://github.com/eternal-november/eos-research-programmes/blob/10e771644cfdf58c3d05e989e4b26cae7164c80a/projects/contextual-composition/protocols/plural-semantics/Plural_Semantics_for_Result_Dependent_Protocols.tex\#L504}{TeX}]} and observation comparison~\mbox{[\href{https://github.com/eternal-november/eos-research-programmes/blob/10e771644cfdf58c3d05e989e4b26cae7164c80a/projects/contextual-composition/protocols/plural-semantics/Plural_Semantics_for_Result_Dependent_Protocols.pdf}{PDF}, \href{https://github.com/eternal-november/eos-research-programmes/blob/10e771644cfdf58c3d05e989e4b26cae7164c80a/projects/contextual-composition/protocols/plural-semantics/Plural_Semantics_for_Result_Dependent_Protocols.tex\#L596}{TeX}]}
make the point precise; the casebook works through its use \cite[Sec.~2]{Casebook}.
Reassociation preserves the causal tree and may change the chronological
word in which the evaluator records events.

\subsection{Replacement needs a relation at every actual continuation}
The finite contextual comparison has a precise scope. Fix the public event
signature, common ranked code library and permitted observations. Interfaces carry
actual capabilities, with their identities and kinds; pure data operations
act only on public data and carry private entries unchanged with their owners.
The common library has identical public tags, captures and templates, though
its bodies may use the compared imported methods. Data operations do not
duplicate ownership. Intermediate leaves retain full private stores.
Choose relations on those stores pairing equal public configurations, and
an initial coupling supported by those relations. That coupling may contain
arbitrary correlations. Each primitive method, and every fixed client
primitive, must supply a correct-marginal finite coupling at every related
actual input and common public argument. Its outputs must have equal
normalized public shapes, related full leaves and nonincreasing candidate
cost. Optional erasure of private unary management uses a fixed cartesian,
leaf-preserving monad map.

Under these hypotheses, the replacement theorem extends the coupling through
all finite well-typed closing contexts of the fixed language, including
finite dependent substitution and acyclic procedure plugging
(see the replacement theorem~\mbox{[\href{https://github.com/eternal-november/eos-research-programmes/blob/10e771644cfdf58c3d05e989e4b26cae7164c80a/projects/contextual-composition/protocols/refinement-recovery/Contextual_Refinement_and_Controlled_Recovery.pdf}{PDF}, \href{https://github.com/eternal-november/eos-research-programmes/blob/10e771644cfdf58c3d05e989e4b26cae7164c80a/projects/contextual-composition/protocols/refinement-recovery/Contextual_Refinement_and_Controlled_Recovery.tex\#L358}{TeX}]} and the worked application
in \cite[Sec.~2]{Casebook}). The support constraint and correct marginals
are the usual ingredients of relational proofs by probabilistic coupling
\cite{barthe}. Here the support condition also records the private stores
needed for the next method calls. Child couplings are multiplied conditional
on the coupled prefix at the actual matched leaves. Under the resulting coupling,
public causal trees agree and candidate cost is no greater than reference cost. If the local costs agree
exactly, the joint observation laws agree. This is neither full abstraction
nor a theorem for arbitrary auditors, schedulers or timing observations.
A newly admitted observation primitive must satisfy the same premises.
The shared latent-state calculation in Appendix~\ref{app:diagnostic}
shows concretely why forgetting private state before filling the holes can
invalidate the multiplication.

Even the joint law \eqref{eq:joint-main} need not preserve an all-raw
resource budget when zero-mass rows are admitted. Adding a zero-probability,
high-cost row leaves that law unchanged but changes which raw executions
satisfy the budget. Cost-lowering native maps also need not preserve exact
cost laws. The separate
law-and-possibility construction (source labels: \texttt{thm:obstruction}, \texttt{thm:K}, \texttt{thm:interpretation})~\mbox{[\href{https://github.com/eternal-november/eos-research-programmes/blob/10e771644cfdf58c3d05e989e4b26cae7164c80a/projects/contextual-composition/protocols/law-and-possibility/Law_and_Possibility_Repair_for_Dependent_Protocols.pdf}{PDF}, \href{https://github.com/eternal-november/eos-research-programmes/blob/10e771644cfdf58c3d05e989e4b26cae7164c80a/projects/contextual-composition/protocols/law-and-possibility/Law_and_Possibility_Repair_for_Dependent_Protocols.tex}{TeX}]}
repairs the former loss for its stated observation interface while still
forgetting witness multiplicity and maps. This is a controlled forgetting
operation, not an equivalence with the native calculus.

\subsection{Robot undo under interference}
An acquired lock can return a journal and enable an undo, while a busy
result returns neither. Those result-dependent interfaces are composed by
the same finite addressed-tree construction. Whether undo preserves other
actors' changes is a further semantic question. For invertible robot steps
$a_k$, let $P_k=a_k\cdots a_1$ be the prefix, and let arbitrary outside
functions $e_k$ act after step $k$. Compensation by $P_m^{-1}$ restores the
state that those same labelled outside functions would have produced on the
untouched initial state, for every initial state, precisely when
\begin{equation}\label{eq:robot-main}
 (P_m^{-1}e_mP_m)\cdots(P_1^{-1}e_1P_1)=e_m\cdots e_1.
\end{equation}
If allowed outside choices include identity and are independent at each
position, validity for every such history is equivalent to commutation of
each allowed $e_k$ with its corresponding $P_k$. A reset need not commute,
even when a translation does.

Atomic-looking behavior requires more: admitted observations must be
invariant under the visited prefixes, and the scheduler must make its choices
from the common admitted information. A camera that sees an intermediate
rotation defeats that comparison even when compensation gets the final
state right. Appendix~\ref{app:robot} proves the finite fixture calculation
and states its observation projection. This is a downstream application of
the common sequential calculus together with a specialized interference
argument, not evidence that a generic crossing solves physical recovery.

\section{The boundaries are part of the unification}\label{sec:limits}
The example in Section~\ref{sec:mixed-main} concerns conservativity for a
chosen lawful crossing. There is also an obstruction to unrestricted
conservative amalgamation over a common ordinary category, even before
requiring a presentation by contextads and contentads \cite[Sec.~10]{U}.
The following bicyclic example is stated at that common-base level.

\begin{proposition}[No universal conservative assembly]\label{prop:bicyclic-main}
There are faithful functors from one ordinary category into two others for
which no common target, identifying the pure maps and preserving the given
equations, receives both faithfully. For the associated locally discrete
bicategories, a common bicategorical target cannot receive both locally
fully faithfully, even with coherent pure agreement.
\end{proposition}
\begin{proof}
Take one-object categories
$C=\langle a\rangle$,
$M=\langle a,b\mid ab=1\rangle$,
$N=\langle a,c\mid ca=1\rangle$.
In $M$, represent $a(n)=\max(n-1,0)$ and $b(n)=n+1$ on natural numbers.
The powers of $a$ are distinct and $ba\ne1$. In $N$, use successor for
$a$ and truncated predecessor for $c$; again the pure powers are distinct.
In any common target,
$c=c(ab)=(ca)b=b$, hence $ba=ca=1$. This collapses distinct arrows of
$M$. In a bicategorical target the same calculation gives an invertible
2-cell $ba\cong1$, but the discrete source hom has no such cell. Local
fullness therefore fails.
\end{proof}
This does not assert that these monoids arise from the particular dependent
comma presentations above. It rules out an unqualified conservative
common-base assembly principle. Nor does it prohibit a nonconservative
output. The positive local-full-faithfulness result for probability/list
is consequently substantive, not a consequence of sharing an ambient
category.

The scope of representation is equally important. The project has a
machine-checked obstruction to a discrete graded presentation of the Boolean
continuation monad with tag-preserving pure-postcomposition naturality
\cite{DiscreteObstruction}. It does not exclude all nondiscrete gradings.
Likewise, marked recognition of a crossing in Appendix~\ref{app:recognition}
forces exactness only on displays realized by the actual context calculus.
It is not an intrinsic classification of arbitrary mixed laws.

\subsection{Additional operations require additional structure}
The sequential calculus neither assumes nor automatically supplies a total
feedback operation, concurrency interchange, differentiation or a logical
doctrine. These are different mathematical interfaces. Forward Para models
parameterized composition; derivatives require differential structure and
compatible reparameterization. A stochastic model supplies randomness but
not Bayesian conditioning. A logical interpretation must prove that its
predicates or predicate transformers respect the relevant substitutions.
No general differential or logical enrichment is established here.

There are stronger scoped feedback developments in the project, which should
not be mistaken for either absence of progress or a theorem about every
interactad. The stage- and export-preserving barrier construction gives
synchronized crossing and certified birth-reactive feedback on infinite
prefixes, after a completion-sensitive executed-prefix interface fails its
crossing test. The nonblocking recurrent construction permits repeated paid
work during pending initialization, using finite registries, complete
selected laws and retained historical ownership effects. These are named
changed-interface constructions with their own admission and compatibility
conditions; global admission, centralized control of whole laws and
decentralized execution remain distinct. Appendix~\ref{app:reuse} links the
precise reading routes. Neither supplies arbitrary feedback or unrestricted
concurrency for the present sequential package.

\section{Conclusion}\label{sec:conclusion}
The common calculus is a relative algebra of computations with a visible pure
interface, native maps and changing-endpoint squares. Its context and content
presentations explain different sides of a program's interface. Their
normalization gives one category of mixed computations; the crossing explains
how the kinds of computation act on one another during composition. Ordinary
fixed effects, forward parameters, display witnesses and produced content
therefore enter through specific recoveries and comparison maps rather than
through a list of analogies.

The probability/list construction demonstrates the central interaction in
full: an outcome changes the addressed family of subsequent requirements,
and those requirements contribute their actual witnesses and probability
maps to the composite. The independent-input theorem provides a reusable
way to construct such laws, while the receipt example shows why coefficient
and content cells must both survive. Observation theorems then justify
particular implementations or abstractions at a stated interface. They do
not redefine equality of native computations, and compensation theorems do
not follow from composition alone.

Further work should preserve that separation of obligations. A new effect
combination needs a lawful crossing; a conservative combination needs a
hom-category recovery proof; a coarser semantics needs a composition and
observation comparison; a feedback extension needs its own admissible
closure and native compatibility equations.
The aim is a common calculus whose useful specializations can be combined
and audited without silently forgetting what made them different.

\appendix
\section*{Appendices}
Appendix~\ref{app:credit} compares the construction mechanisms with prior work. Appendix~\ref{app:construction} states and proves the independent-input theorem and its Set realization. The indexed and recognition examples follow, then the full local equations, pasting checks and compatible input changes. The diagnostic and robot calculations specify their own observation and interference assumptions. Appendix~\ref{app:reuse} records the source and verification route.

\section{Construction mechanisms and prior work}\label{app:credit}
The wreath formalism already provides a place for a coherent crossing:
Capucci--Myers include an intertwiner and its structure cells among the
wreath data \cite[Def.~D.9]{CM}. Theorem~\ref{thm:glue} constructs the intertwiner and its structure cells
from the two input models. Two
established construction mechanisms are particularly close. Buckley--Garner's
co-crossed double categories have codomain-cartesian lifts stable under
horizontal composition, and their Proposition~6.3 constructs a crossing
from them \cite[Def.~6.1, Prop.~6.3 and Rem.~9.5]{BG}. The relative
construction uses that cartesian-square mechanism along context displays.
Its universal property tests every native content square with the
specified right edge. All pure maps remain in the base, and descent to
the relative tensor requires a separate cartesian comparison for pure
balance.

Cnossen, Haugseng, Lenz and Linskens identify commutative semiring objects
in cartesian closed presentable $\infty$-categories with product-preserving
finite-bispan functors \cite[Thm.~A]{CHLL}.
The raw category-valued coefficient calculation below belongs to that
calculus. Directed additive transfer and the slice imposing a total
augmentation supply further data; the augmented category itself need not
be a product-preserving bispan model. The contribution is the construction
that matches this decorated occurrence calculus to independently supplied
relative content corners and descends through arbitrary pure balance.
A reduction of the whole construction to a suitably decorated
corners-and-bispans theorem remains open.

\paragraph{Relative polynomial geometry.}
Von Glehn constructs the law of free fibred products over free fibred
sums and treats a class closed under composition and arbitrary base
change \cite[Thms.~7.1--7.2 and Sec.~8]{vonglehn}. Thus the use of a
restricted, pullback-stable class is not new. Weber composes polynomials
with exponentiable middle arrows over categories with pullbacks by
distributivity pullbacks, and interprets them by pullback, dependent
product and dependent sum on slices
\cite[Thms.~3.1.10 and 3.2.6]{weber}. Here the operands are given
content and context calculi, not free fibred completions. A rig tensor
is not in general a dependent product, as the dimension calculation
in Section~\ref{sec:paired} illustrates. Gambino and Kock's polynomial
substitution and cartesian position squares give the bare addressed
geometry \cite[Secs.~1.11, 2.1 and 4.1]{gambinokock}; the map-level
work here is to identify that geometry over the original grade and
coefficient maps and prove compatibility with pure balance.

\paragraph{Indexed transfers and representation.}
Ponto and Shulman compose using external tensor, diagonal restriction
and indexed coproducts, which are left adjoints to restriction
\cite[Def.~3.1 and Thm.~5.2]{pontos}. Elmanto and Haugseng's bispan
extension uses left adjointability and invertible distributivity mates
\cite[Def.~2.4.10, Thm.~2.5.1 and Cor.~2.5.15]{eh}. Chosen directed
additive transfer here need not satisfy that adjoint prescription:
in finite-dimensional vector spaces with only surjective linear maps,
the fold transfer sends $(k,k)$ to $k^2$. It has an identity map to
$k^2$, whereas there is no map from $(k,k)$ to the constant family
$(k^2,k^2)$. The failure concerns that transfer, not the scope of
bispan theories with additional structure. The double Grothendieck
representation of Cruttwell, Lambert, Pronk and Szyld
\cite[Def.~2.23 and Thm.~3.41]{doublefib} gives a different organization
where full double-fibration hypotheses hold. It is not invoked for
the codomain fibration alone, and does not substitute for checking
the common occurrence lift or its pure balance.

\paragraph{Constructing a law versus using one.}
Street's wreath products and mixed Kleisli convolution illustrate what
specified interaction cells can build \cite[Secs.~2--3]{street};
Marmolejo's Gray-category theory relates supplied pseudodistributive
laws, compatible composites and liftings \cite{marmolejo}. These are
formal precedents for the composition step, not the construction of
the interaction cells in this paper. There are also genuinely
constructive Set precedents: Manes and Mulry build ordinary laws using
recursive constructions and quotients by linear equations \cite{manes}.
Fujii, Tsai, Montacute and Hasuo give a canonical Set construction from
a $W$-operadic monad and a $W$-commutative monad
\cite[Thm.~42 and Rem.~43]{fujii}. Their full monad--monad law uses an
additional condition: ordinary commutativity of the second monad or
absence of operadic operations of arity greater than one. Finite
distributions satisfy the commutativity condition. With lists, the
resulting Set law is the output-law comparison of Section~\ref{sec:prob}.
The native law of Theorem~\ref{thm:problaw} composes positive finite
witness spaces and their deterministic measure-preserving maps. The
separate coefficient and grade cells of the general construction are
developed in Appendix~\ref{app:construction}.

The independent-input theorem is a sufficient construction. Appendix~\ref{app:recognition}
asks a different question: which conditions can be recovered from a mixed
law that preserves the original projections? That marked recognition
statement does not classify arbitrary laws, and does not extend its tests
to displays absent from the context calculus. The original full written gluing proof is \cite{Gluing}, with prior-work
analysis in \cite{Comparison}. Appendix~\ref{app:construction} and its
coherence appendices present that construction within the broader
unification of this article. Appendix~\ref{app:reuse} locates the source
and distinguishes written proofs from formalized subresults.

\section{The independent-input construction}\label{app:construction}
The complete construction used in Section~\ref{sec:independent-main} follows. It applies beyond the finite probability/list model and keeps the decorated context maps and native content cells separate.

\subsection{The common interface: occurrences and two pullbacks}\label{sec:interface}
In the two-resource example (Section~\ref{sec:paired}), the receipt
$[*,*]$ requests two continuations with the same label. The list records
two addresses, so each continuation can make its own resource choice.
A cartesian monad $T$ abstracts this occurrence interface. Its unit
describes the ordinary continuation and its multiplication substitutes
continuation shapes. Cartesianity supplies the pullback comparisons
for substitution. The decorated lift below specifies the resource
operation at those addresses.

Fix an ordinary category $\B$ with pullbacks and a class $\J$ containing
isomorphisms, closed under composition, pullback and isomorphism of arrows.
Members of $\J$ will be called \emph{displays}. Assume $T$ preserves
pullbacks, its unit $\eta$ and multiplication $\mu$ are cartesian, and
$T\J\subseteq\J$. All pure maps of $\B$ remain available; the word
``relative'' restricts display lifting, not the pure category. Work in
universes in which the category-valued pseudocoends used below exist.
For arrows, $vu$ means $v\circ u$; words of computations are written in
execution order. For $a:A\to TB$ and $b:B\to TC$ put
$b\bullet a=\mu_C(Tb)a$.

A content arrow $n:A\rightsquigarrow B$ is assigned an occurrence arrow
$a_n:A\to TB$. A decorated display $P:E\to B$ lies over $r:E\to B$.
The desired crossing must have carrier
\begin{equation}\label{eq:carrier}
                         A_r=A\times_{TB}TE.
\end{equation}
There are two different constructions of this carrier. On the resource
side, apply $T$ to the display, with its decoration, and restrict along
$a_n$. On the content side, lift the actual arrow $n$ along $r$. The
second construction must retain its native content cell, not merely its
image under $a$. The theorem matches the constructions by
\eqref{eq:carrier}; it does not identify their maps.

\subsubsection{Decorated displays and their occurrence lift}\label{sec:decinput}
A pseudo-double category $\D$ records contexts. Its vertical category is
$\B$, and a forgetful pseudo-double functor to commuting squares of $\B$
is identity on vertical arrows. A horizontal arrow $P:E\to B$ lies over
a display $r:E\to B$; a square $P\to P'$ lies over
$(h:E\to E',u:B\to B')$ with $r'h=ur$. Distinct decorated squares
may have that same boundary. For $S:F\to E$, write $P\st S:F\to B$
for horizontal composition; it lies over the composite display.
Horizontal units $\un_B$ lie over identities. Associators and unitors
have identity boundary; carrier transport by isomorphisms is understood.

Require $\cod:\D_1\to\B$ to be a fibration, with chosen cartesian
squares whose underlying squares are pullbacks. Require horizontal units
and horizontal pastes of cartesian squares to be cartesian. Thus $u^*P$
exists for \emph{every} pure $u$, and restrictions commute with accumulation
up to their canonical comparisons. In the fibre $\K_B$, arrows are
squares over $\id_B$; the carrier functor is $e_B(P)=E$. The contextful
module is
\[
 M(A,B)=(e_A\downarrow B),\qquad
 (P,g)=\bigl(A\xleftarrow{r}E\xrightarrow{g}B\bigr).
\]
Its composition is restriction followed by $\st$, as above. Cartesian
uniqueness identifies the reindexed displays in triple composites; the
horizontal coherence of $\D$ then proves the associativity and unit laws
on the actual decorated maps.

\begin{definition}[Occurrence structure on contexts]\label{def:lift}
An occurrence lift consists of a pseudo-double endofunctor $\Tt$ over
$T$, preserving codomain-cartesian squares and taking $P$ to a decoration
of $Tr:TE\to TB$, together with vertical double transformations
\[
 \bar\eta:1\Rightarrow\Tt,\qquad
 \bar\mu:\Tt^2\Rightarrow\Tt.
\]
The horizontal comparisons
$\Tt(P\st S)\cong\Tt P\st\Tt S$ and
$\Tt\un_B\cong\un_{TB}$ are invertible and satisfy the functor
pentagon and unit equations, naturally on all squares. The transformations
lie over $\eta,\mu$; their squares on each display are codomain-cartesian.
They are natural on the vertical and horizontal-arrow categories and
commute with horizontal units and composition, inserting these functor
comparisons. Finally, as decorated transformations,
\[
 \bar\mu\,\Tt\bar\mu=\bar\mu\,\bar\mu_{\Tt},\qquad
 \bar\mu\,\Tt\bar\eta=1=\bar\mu\,\bar\eta_{\Tt}.
\]
\end{definition}
This is a condition on contexts alone. In particular the multiplication
axiom does not refer to a content grade. For sets it will follow from
addressed finite polynomial evaluation; it is stated independently of
that presentation so that an indexed example can use the same theorem.

\subsubsection{Content with relative cartesian corners}\label{sec:coninput}
To move a resource display past a content computation, the construction must compare
content arrows even when their input and output interfaces change.
For fixed interfaces \(A,B\), the category \(N(A,B)\) contains
the content arrows \(A\rightsquigarrow B\) and their native comparison
cells. Composition and identities are coherent up to isomorphism.
Together with the actions of pure maps this is a relative
pseudopromonad \(N:\B^{\op}\times\B\to\Cat\). Its realization is a
bicategory on the same objects as \(\B\), with an identity-on-objects
pure pseudofunctor \(\pure:\B\to N\) \cite[Secs.~2--3]{U}.

For a pure input map \(k:A\to A'\) and output map \(u:B\to B'\),
precomposition with \(k\) and postcomposition with \(u\) transport
\(n':A'\rightsquigarrow B'\) and \(n:A\rightsquigarrow B\) into the
same category \(N(A,B')\). A square from \(n\) to \(n'\) over
\((k,u)\) is the native cell \(u_*n\to k^*n'\) in that category.
These squares form the double category \(\N\), using the action
comparisons and the original equality of native cells. The square
\(n\to u_*n\) is opcartesian for right-edge transport: the comparison
\((vu)_*\cong v_*u_*\) gives the unique factor with right edge \(v\). The
construction uses only the native content categories and pure actions.

The occurrence assignment must preserve the whole content calculus:
\begin{equation}\label{eq:arity}
 a_{\pure(u)}=\eta u,\qquad a_{nm}=a_m\bullet a_n,
 \qquad Tu\,a_n=a_{n'}k\quad\text{for every square over }(k,u).
\end{equation}
It is a multiplicative pure-compatible map to the locally discrete
Kleisli category. It need not be faithful on native content maps.

\begin{definition}[Relative content corners]\label{def:corners}
For every content arrow $n:A\rightsquigarrow B$ and display $r:E\to B$
choose a square
\begin{equation}\label{eq:corner}
 C_r(n)=\left[
 \begin{array}{ccc}
 A_r&\xrightarrow{n_r}&E\\
 {\scriptstyle\ell}\downarrow&&\downarrow{\scriptstyle r}\\
 A&\xrightarrow{n}&B
 \end{array}\right].
\end{equation}
The following properties are required.

\smallskip\noindent
\textup{(i)} The square is cartesian for $\cod:\N_1\to\B$: a square
$z\to n$ with right edge $rv$ factors uniquely through $C_r(n)$ with
right edge $v$. Testing squares and their left and right edges are
unrestricted pure maps.

\smallskip\noindent
\textup{(ii)} Its image is the occurrence pullback:
$(\ell,a_{n_r}):A_r\iso A\times_{TB}TE$. Horizontal identities on
displays are cartesian. Matched horizontal pastes of these cartesian
squares are cartesian, lifting the last content arrow first and then its
predecessor along the resulting left edge. That edge belongs to $\J$.

\smallskip\noindent
\textup{(iii)} For every pullback of a display along an arbitrary pure map,
\[
 \begin{array}{ccc}E_u&\xrightarrow v&E\\
 {\scriptstyle r_u}\downarrow&&\downarrow{\scriptstyle r}\\
 B&\xrightarrow u&B',\end{array}
\]
the canonical square $\pure(v)\to\pure(u)$ over $(r_u,r)$ is
codomain-cartesian. It is the pure comparison of $rv=ur_u$, not an
independently chosen square with isomorphic corners.
\end{definition}
These properties have different jobs. Cartesianity factors native content
maps. The occurrence pullback makes their boundaries match decorated
maps. Horizontal stability controls successive content computations.
The pure square is what permits a change of interface to move from one
side of a middle boundary to the other.

\subsection{Constructing the crossing}\label{sec:construction}
The crossing moves a resource display past a content arrow. The two
constructions of its carrier in \eqref{eq:carrier} must match, while
the resource map and the content cell remain separate as in the
two-resource calculation (Section~\ref{sec:paired}). For a module $X$, let $\grade_X(A,C)$ have objects $(P,x)$,
where $P\in\K_A$ and $x\in X(e_AP,C)$. An arrow is a pair
\begin{equation}\label{eq:normalmap}
 (\alpha:P\to P',\quad t:x\to h^*x'),\qquad h=e_A\alpha.
\end{equation}
For endpoint changes $(k,v)$, use a decorated square over $(h,k)$ and
$t:v_*x\to h^*x'$. Equality is tested on $\alpha$ and $t$ separately.
This is not an equality test on $h$.

For modules $X,Y$ use the category-valued profunctor tensor of
Lawler \cite[Sec.~4.1.1, Eq.~(4.1.1)]{Lawler}, restricted to ordinary
base categories and written in execution order:
\[
 (X\odot Y)(A,C)=\int^{B\in\B}_{\mathrm{ps}}X(A,B)\times Y(B,C).
\]
The pseudocoend retains pairs of native arrows and adjoins invertible
balances $[u_*x,y]\cong[x,u^*y]$ for \emph{every} pure $u$, with
identity, composition and naturality relations. Normalization gives
\begin{equation}\label{eq:normal}
 M\odot X\simeq\grade_X,
 \qquad [B;(P,g),x]\longmapsto(P,g^*x).
\end{equation}
Indeed a reverse functor chooses $B=e_AP$ and $g=\id$. On
\eqref{eq:normalmap}, apply $t$, balance along $h$, then apply $\alpha$.
The balance along $g$ compares an arbitrary representative to this one;
the balance relations give the two inverse comparisons. This also proves
naturality at the endpoints. In particular the tensor is not a set of
isomorphism classes of diagrams.

On contexts define the occurrence action
\begin{equation}\label{eq:L}
 L_aP=a^*\Tt P,
 \qquad a_P:A\times_{TB}TE\longrightarrow TE
 \quad(a:A\to TB, P:E\to B).
\end{equation}
The second projection $a_P$ is itself a Kleisli arrow to $E$. The lifted
monad supplies four useful comparisons:
\begin{align}
 I_a &:L_a\un_B\cong\un_A,&
 U_u &:L_{\eta u}P\cong u^*P,\label{eq:IU}\\
 K_{a,b}&:L_aL_bP\cong L_{b\bullet a}P,&
 F_a&:L_a(P\st S)\cong L_aP\st L_{a_P}S.\label{eq:KF}
\end{align}
Here $K$ uses the cartesian multiplication square of $\bar\mu$,
$U$ the unit square of $\bar\eta$, and $F,I$ the horizontal functor
comparisons. Their full equations are proved in Appendix~\ref{app:local}.
The important mixed equation says that flattening occurrences and
accumulating contexts commute on decorated maps:
\begin{align}
 F_{b\bullet a}(P,S)K_{a,b}(P\st S)
 ={}&\bigl(K_{a,b}(P)\st K_{a_{L_bP},b_P}(S)\bigr)\notag\\
 &\quad\cdot F_a(L_bP,L_{b_P}S)L_a(F_b(P,S)).\label{eq:middle}
\end{align}
The second slot is transported along the carrier comparison of $K$.
Its typing equation is
$b_P\bullet a_{L_bP}=(b\bullet a)_P\dom K_{a,b}(P)$.
This is the single place where the two kinds of accumulation meet on
the context side. It follows from the horizontal-composition axiom of
$\bar\mu$, not from an assumed mixed law.

\begin{theorem}[Composition from independent inputs]\label{thm:glue}
Assume the base, display class and cartesian monad of
Section~\ref{sec:interface}, the decorated display calculus with the
occurrence structure of Definition~\ref{def:lift}, and the content
calculus satisfying \eqref{eq:arity} and Definition~\ref{def:corners}.
There is a pseudodistributive law
\[
                    \tau:N\odot M\longrightarrow M\odot N
\]
whose normalized value on $[B;n,(P,g:E\to C)]$ is
\begin{equation}\label{eq:cross}
                         (L_{a_n}P,\ g_*n_r).
\end{equation}
The source of $n_r$ is transported by the specified pullback comparison
$A_r\cong A\times_{TB}TE$. The law acts on the native maps, is balanced
over the whole pure category $\B$, and has the four comparisons specified
below. Consequently $\grade_N$ has the composite relative pseudopromonad
structure, with multiplication obtained by the middle crossing
\[
 (M\odot N)\odot(M\odot N)
 \xrightarrow{\,M\odot\tau\odot N\,}
 M\odot M\odot N\odot N
 \xrightarrow{\,m_M\odot m_N\,}M\odot N.
\]
The construction is natural under the compatible maps of independent
inputs specified in Appendix~\ref{app:parameters}.
\end{theorem}

\subsubsection{Why the two native maps fit}
Take a content square $s:n\to n'$ over $(k,u)$ and a context square
$P\to P'$ over $(h,u)$. Cartesianity produces a unique content square
$\widetilde s$ satisfying
\begin{equation}\label{eq:factor}
 C_{r'}(n')\widetilde s=s C_r(n),
 \qquad\text{right edge of }\widetilde s=h.
\end{equation}
If $H$ is its left edge, \eqref{eq:arity} says that its coordinates in
the target pullback are $(k\ell,Th\,a_{n_r})$. Applying $\Tt$ to the
given decorated square and then restricting along $a_n$ gives a decorated
square with exactly those coordinates. Therefore it and $\widetilde s$
form an arrow of \eqref{eq:normalmap}. With tails satisfying $vg=g'h$,
the content component is the bent cell
\[
 v_*g_*n_r\cong g'_*h_*n_r
       \longrightarrow g'_*H^*n'_{r'}\cong H^*g'_*n'_{r'}.
\]
The middle arrow is not required to be invertible. Identity and composition
follow by restoring $C_{r'}(n')$ in \eqref{eq:factor} and using its unique
factorization, while the decorated factors compose in $\D$. Thus the
construction already retains both kinds of maps in the receipt example.
For a left endpoint change, paste the pure pullback square for that change
with $C_r(n)$; horizontal stability identifies the result with the lift
of the restricted content arrow. This gives the invertible endpoint
comparison. A right endpoint change is simply transport of the final tail.

There is a local way to describe precisely what has been matched. Put a
superscript $0$ on $\grade$ when displays are undecorated and let
$K_T(A,C)=\Disc\B(A,TC)$. Then
\begin{equation}\label{eq:pullcat}
 \begin{array}{ccc}
 \grade_N^{\D}&\longrightarrow&\grade_N^0\\
 \downarrow&&\downarrow\\
 \grade_{K_T}^{\D}&\longrightarrow&\grade_{K_T}^0
 \end{array}
\end{equation}
is a pullback of categories in a fixed named-carrier presentation.
A matched object is exactly $(P,n)$; a matched arrow consists of the
actual decorated arrow and the actual content cell with the same carrier
map. Composition is the componentwise composition already used above.
It is also a pseudopullback up to equivalence: the right vertical functor
is an isofibration, since an isomorphism of source carriers lifts by pure
precomposition with its inverse. An isomorphic matching can therefore be
rectified. This local property explains the gluing of maps; it is not an
initiality property for all possible mixed calculi.

\subsubsection{Why arbitrary pure changes of interface are allowed}
Suppose $u:B\to B'$ is pure and $P$ lies over $r:E\to B'$. Pull back
$r$ to $r_u:E_u\to B$, with $v:E_u\to E$. Crossing after $u_*n$ and
crossing $n$ against $u^*P$ must represent the same operation. On content,
the required comparison is the canonical cartesian--opcartesian mate
\begin{equation}\label{eq:mate}
                    v_*n_{r_u}\iso H^*(u_*n)_r.
\end{equation}
This is the usual Beck--Chevalley comparison of cartesian and
opcartesian completions \cite[Def.~13.21 and Lem.~16.1]{shulman}.
The carrier $H$ is invertible because $T$ preserves the pullback and
$a_{u_*n}=Tu\,a_n$. That fact alone does not prove the cell invertible.
For the cell, paste $C_{r_u}(n)$ horizontally with the \emph{canonical}
pure pullback square $\pure(v)\to\pure(u)$. Definition~\ref{def:corners}
makes this paste cartesian. Its bottom is $u_*n$, its right edge is $r$,
and its top is $v_*n_{r_u}$. Comparison to $C_r(u_*n)$ by cartesian
uniqueness is exactly \eqref{eq:mate}. On decoration, the corresponding
map is
\[
 a_n^*\Tt(u^*P)\cong a_n^*(Tu)^*\Tt P
                  \cong(Tu\,a_n)^*\Tt P.
\]
It has the same carrier $H$. Thus the two components supply the balance,
including the original content mate rather than an unspecified isomorphism.

For a composite of pure maps, restore the final cartesian arrow and
unbend the pure transports. The two proposed mates then have the same
pasted pure pullback and the same initial square. Cartesian uniqueness
makes them equal. This proves composition of balances; the identity case
and their naturality on \eqref{eq:factor} follow by the same cancellation.
The decoration calculation is cartesian uniqueness in $\D_1$.
These are precisely the defining pseudocoend relations, so the crossing
descends. Conversely, putting $n$ equal to a horizontal identity in
\eqref{eq:mate} shows that invertibility of all these canonical mates
recovers pure exactness. The pure hypothesis is not redundant bookkeeping.

\subsubsection{The four comparisons and the composition proof}
The law must commute with each of the two units and each of the two
multiplications. Let $-P$ denote the context with identity forward tail.
For consecutive content arrows $n,m$ and displays $P,S$, the comparisons
are obtained as follows; inverses are inserted where the displayed
orientation requires them.
\begin{center}\small
\begin{tabular}{@{}p{.14\linewidth}p{.28\linewidth}p{.51\linewidth}@{}}
\toprule
Operation&Context component&Content component\\\midrule
Content unit&$U_{\id}(P)$&Compare $C_r(1)$ with the horizontal identity on $r$.\\[3pt]
Context unit&$I_{a_n}$&Compare $C_{\id}(n)$ with the identity square on $n$.\\[3pt]
Content multiplication&$K_{a_n,a_m}(P)^{-1}$&Compare $C_r(nm)$ with the horizontal paste of $C_\ell(n)$ and $C_r(m)$.\\[3pt]
Context multiplication&$F_{a_n}(P,S)$&Compare $C_{rs}(n)$ with the vertical paste of $C_s(n_r)$ and $C_r(n)$.\\\bottomrule
\end{tabular}
\end{center}
Here $\ell$ is the left edge of $C_r(m)$ in the third row. Each pair of
content squares is cartesian with the same bottom and right edge, up to
the one-sided structural comparisons. Hence there is a unique invertible
comparison of their tops. Its occurrence image is forced by the two
pullback coordinates. In the third row these are the original source
coordinate and $\mu$-flattening; in the fourth they are the source
coordinate and the final $TF$-coordinate. The units use the cartesian
unit of $T$ and the identity pullback. These are exactly the carriers
of the context comparisons. Each row is therefore one arrow of
$\grade_N$, not two unrelated isomorphisms.

Their types are
\begin{align}
 \tau(j_NM)&\cong Mj_N,&\tau(Nj_M)&\cong j_MN,\notag\\
 \tau(m_NM)&\cong(Mm_N)(\tau N)(N\tau),&
 \tau(Nm_M)&\cong(m_MN)(M\tau)(\tau M).\label{eq:four}
\end{align}
To prove coherence, consider a rectangle with content columns and display
rows. Fill it from the lower right with chosen cartesian squares. Every
subrectangle is cartesian, by horizontal stability and ordinary
composition of cartesian arrows. Any sequence of comparisons of
subrectangles, followed by comparison to the outside rectangle, has the
same restored content square. The one-sided associators and unitors give
the same restored boundary by their pentagon and triangle laws. Cartesian
uniqueness therefore proves equality of the content components. This
argument works even though occurrence does not detect equality of content
cells.

The context components require a different argument: they are the $K$,
$F$, $I$, $U$ equations, in particular \eqref{eq:middle}. Appendix~\ref{app:coherence}
records the ten pseudodistributive pastings and matches each to its
context equation and content rectangle. It also checks the internal pure
seams, so that the four comparisons really are modifications on relative
tensors rather than merely operations on identity-tail representatives.
Together those checks prove Theorem~\ref{thm:glue}. The general composition
of pseudomonads with a pseudodistributive law then gives its stated
multiplication \cite[Def.~4.1 and App.~A]{GL}.

\subsection{Resources: polynomial evaluation with directed maps}\label{sec:resources}
This section builds the resource input of Theorem~\ref{thm:glue}:
the decorated occurrence lift of Definition~\ref{def:lift}. A finite
polynomial shape indexes a tensor of resource choices, and directed
transfer gives its action on the merges from Section~\ref{sec:paired}.
Let $\B=\Set$ and let
\[
                    TX=\coprod_{s\in S}X^{F_s}
\]
be a cartesian polynomial monad with finite, addressed position sets $F_s$.
Let $\V$ be a symmetric rig category: $(\oplus,0)$ and $(\otimes,I)$
are symmetric monoidal, with coherent invertible finite distribution and
absorption maps. Neither operation is assumed to be a categorical product
or coproduct. Finite indexed brackets and their canonical reindexings are
fixed. Context displays have finite fibres and belong to $\J$; if the
ambient $\J$ is larger, this finite-fibre subclass is the one used here.

Over $r:E\to B$, choose objects $V_e\in\V$ and, for each $b$, one
augmentation
\begin{equation}\label{eq:augmentation}
                  \epsilon_b:\bigoplus_{e\in E_b}V_e\longrightarrow I.
\end{equation}
A global augmentation is part of the data. Without coproduct inclusions
it cannot be replaced by an assumed family of maps $V_e\to I$.
A map over a square $(h,u)$ consists of the arrows
\begin{equation}\label{eq:aggregation}
 \gamma_{b,e'}:\bigoplus_{\substack{e\in E_b\\he=e'}}V_e
                         \longrightarrow V'_{e'}
       \quad(e'\in E'_{ub}),
\end{equation}
whose sum, after repartitioning the source fibre, preserves
\eqref{eq:augmentation}. Even if $u$ merges base points, the source sum
is over the specified fibre at $b$. Composition partitions the fibre of
a composite $h'h$ into fibres of $h$ and $h'$, then applies the given
arrows. This is why the boundary function does not determine the map.
Restrictions along $u$ use ordinary pullbacks and the same coefficient
objects and arrows under fibre bijections. Their factorization property
against all squares is immediate from that formula.

\subsubsection{What finite bispans compute, and what must be added}
The raw operation associated to a finite bispan is
\begin{equation}\label{eq:bispan}
 \left(L\xleftarrow{s}U\xrightarrow{p}D\xrightarrow{t}K\right)
 \longmapsto
 \left[(V_l)\longmapsto
   \left(\bigoplus_{td=k}\ \bigotimes_{pi=d}V_{si}\right)_{k\in K}\right].
\end{equation}
This is category-valued: tuples of actual arrows of $\V$ are evaluated,
not just their objects. Substitution distributes products over sums;
a resulting term records one alternative at each input position and all
its addressed leaves. A three-level substitution has the same such
histories on both associations, giving the compositor by rig coherence.
This is the finite-bispan coefficient calculus described abstractly by
\cite[Thm.~A]{CHLL}; on an infinite base it is used fibre by fibre.

To obtain \eqref{eq:aggregation}, use the covariant additive transfer
$h_!V=(\bigoplus_{he=e'}V_e)_{e'}$ and its Grothendieck category $\mathcal R$.
An arrow of $\mathcal R$ is $(h,\gamma:h_!V\to V')$. Total sum defines
$\Sigma:\mathcal R\to\V$ by fibre partition. The augmented category is
\[
                 \mathcal R\times_{\V}\V_{/I},
\]
where $\V_{/I}\to\V$ takes the domain. This constructs the directed
maps and their augmentation equation; it does not infer directed maps
between carrier shapes from invertible bispan reindexings. Nor is the
slice generally another raw bispan model. In the discrete rig
$(\mathbb N,+,\cdot)$, an augmented singleton has one possibility, while
an augmented two-point carrier has the two possibilities $(1,0),(0,1)$.
Their fibres therefore do not obey the disjoint-union product rule of
the unaugmented model.

\subsubsection{Accumulating resources before merging them}
For $P:E\to B$ with weights $V_e$ and $S:F\to E$ with weights $W_z$,
the row of $P\st S$ at $z$ is $V_{sz}\otimes W_z$. Its augmentation
first distributes by $e$, applies each inner augmentation, then applies
the outer one. The horizontal unit has one row $(I,\id_I)$.

The map on a composite is worth writing down, because an apparent
shortcut would be wrong. Suppose the outer map has carrier $h$ and
coefficients $\alpha$, and the inner map has carrier $k$ and coefficients
$\beta$. Write $u$ for the outer base map. Fix $b\in B$, a target
$e'\in E'_{u(b)}$, and a target $z'$ with $s'z'=e'$. In every outer
sum below $I_{b,e'}=\{e\in E_b\mid he=e'\}$ is taken in the single
source fibre at $b$. The coefficient map is
\begin{align}
 \bigoplus_{e\in I_{b,e'}}\ \bigoplus_{\substack{sz=e\\kz=z'}}V_e\otimes W_z
 &\cong\bigoplus_{e\in I_{b,e'}}V_e\otimes
                \left(\bigoplus_{\substack{sz=e\\kz=z'}}W_z\right)
 \xrightarrow{\oplus_e(1\otimes\beta_{e,z'})}
              \bigoplus_{e\in I_{b,e'}}V_e\otimes W'_{z'}\notag\\
 &\cong\left(\bigoplus_{e\in I_{b,e'}}V_e\right)\otimes W'_{z'}
 \xrightarrow{\alpha_{b,e'}\otimes1}V'_{e'}\otimes W'_{z'}.
 \label{eq:dependent}
\end{align}
The maps $\beta_{e,z'}$ can depend on the ancestor $e$. They must be
applied before merging those ancestors. Only their common target
$W'_{z'}$ can be taken outside the merge.

This formula proves horizontal functoriality on arbitrary squares.
For two stages of maps, partition by the intermediate ancestor and row.
The second inner map is constant on each already merged intermediate
ancestor, so tensor bifunctoriality moves it past the first outer map.
After distribution this is exactly \eqref{eq:dependent} for the composite.
The same calculation, with the augmentations in place of the target
maps, proves preservation of total augmentation. Triple accumulation
uses tensor associativity on each addressed chain; both augmentations
contract that chain from the inside out. Its naturality follows from the
same constant-on-the-merged-fibre exchange. Tensor coherence supplies
the pentagon and units. Thus the formula defines the required
pseudo-double category, not merely a multiplication on its objects.
Horizontal pastes of cartesian restrictions remain cartesian: their
fibre maps are bijective and their row comparisons are isomorphisms.

\subsubsection{Applying an occurrence shape}
For $t=(s,(b_i)_{i\in F_s})\in TB$, the fibre of $Tr:TE\to TB$ is
$\prod_iE_{b_i}$. Define its row at a choice tuple $(e_i)$ to be
$\bigotimes_iV_{e_i}$ and its augmentation to be
\begin{equation}\label{eq:occaug}
 \bigoplus_{(e_i)}\bigotimes_iV_{e_i}
 \cong\bigotimes_i\bigoplus_{e_i}V_{e_i}
 \xrightarrow{\otimes_i\epsilon_{b_i}}\bigotimes_i I\cong I.
\end{equation}
This defines $\Tt$ on objects. Its action on the native map
\eqref{eq:aggregation}, at a target tuple $(e'_i)$, is
\begin{equation}\label{eq:occmap}
 \bigoplus_{\forall i\; (he_i=e'_i)}\bigotimes_iV_{e_i}
 \cong\bigotimes_i\bigoplus_{he_i=e'_i}V_{e_i}
 \xrightarrow{\otimes_i\gamma_{b_i,e'_i}}\bigotimes_iV'_{e'_i}.
\end{equation}
Each sum remains in the designated $b_i$-fibre. Formula
\eqref{eq:rankmaps} is its two-position instance.

For composites of maps, partition the source tuples by their intermediate
tuple. Distributing this partition position by position gives the same
tensor of coefficient composites as applying \eqref{eq:occmap} twice.
Naturality of the distributor proves equality; totalling the formula and
using the incoming augmentation equations proves \eqref{eq:occaug} is
preserved. Pullback restrictions are preserved because their fibre
bijections and invertible row maps remain so under finite tensor.

\phantomsection\label{par:dependent-fubini-maps}
\noindent\textbf{Dependent Fubini on maps.}
The horizontal comparison splits a tuple of dependent pairs into the
outer tuple and its dependent inner tuple. On coefficients it is the
specified shuffle
\[
 \bigotimes_i(V_i\otimes W_i)
       \cong\left(\bigotimes_iV_i\right)\otimes
                   \left(\bigotimes_iW_i\right).
\]
To check its naturality, expand \eqref{eq:dependent} at every position in
\eqref{eq:occmap}. Both routes sum over the same compatible pairs
$(e_i,z_i)$ and apply $\beta_{e_i,z'_i}$ before $\alpha_{b_i,e'_i}$.
All remaining exchanges are tensor bifunctoriality on disjoint factors
or distributor naturality with a fixed target row. The augmentations
are likewise inner-then-outer on either route. Three displays give the
same block permutation of the same ancestor chains, proving the
horizontal functor pentagon. The unit comparison removes the forced
$I$ factors.

Cartesianity of the polynomial unit gives one position; cartesianity
of multiplication gives a bijection
\[
 F_{\operatorname{flat}(s,(s_i))}\cong\coprod_iF_{s_i}.
\]
These statements follow by evaluating the polynomial transformations
on their generic labellings. Define $\bar\eta$ by the singleton row
identification and $\bar\mu$ by addressed tensor flattening. Their
underlying squares are pullbacks and their row maps are invertible,
so they are cartesian. Naturality on \eqref{eq:occmap} is the same tensor
of leaf aggregations on both sides. For horizontal compatibility of
$\bar\mu$, both routes collect all outer row factors and all inner row
factors in the same leaf order. The monad equations identify the leaf
bijections, not just the number of leaves; tensor coherence completes
the equations on decorated arrows. This proves every condition in
Definition~\ref{def:lift} without referring to a content grade.

These coherence arguments use addressed occurrences. Independent factors
are tagged before equal values are substituted; histories and their leaf
bijections determine the structural maps. Regular distributive coherence
\cite{Laplaza} applies with those tags, including distinct unit variables
for nullary nonzero alternatives before specialization to $I$.
Nonstructural arrows move through structural maps only by naturality
where incidences agree. Empty sums and products use absorption and unit
coherence, not assumed uniqueness of arrows from $0$ or to $I$.

\subsection{Content: exact recovery of the comma presentation}\label{sec:recovery}
Receipt grades in Section~\ref{sec:paired} extend the result by tags
carrying occurrence lists. The following test treats arbitrary comma
content of that form of presentation: it constructs the content corners
of Definition~\ref{def:corners} from maps of grades. Combined with the
resource lift of Section~\ref{sec:resources}, it recovers the original Set
law on those maps. Suppose $p:\Lcat\to\B$ is a global opfibration of grades with extension
$d:\Lcat\to\B$, giving a valid one-sided comma content calculus
\[
 N(A,B)=(A\downarrow d_B).
\]
An object is $(Q,f:A\to dQ)$; a total grade arrow $m:Q\to Q'$ over
$u$ gives a square over $(k,u)$ when $dm\,f=f'k$. Denote the chosen
opcartesian arrow by $\kappa_{u,Q}:Q\to u_!Q$.
The unit is a grade $U_B$ with body $\eta^L_B:B\to dU_B$.
Dependent accumulation of $Q\in\Lcat_B$ and $R\in\Lcat_{dQ}$ is
$Q\di R\in\Lcat_B$, with body
$\mu^L_{Q,R}:dR\to d(Q\di R)$. These bodies need not be invertible.
In particular composition of $(Q,f)$ then $(R,b)$ has grade
$R\di b_!Q$ and body
$\mu^L_{R,b_!Q}\,d\kappa_{b,Q}\,f$.

The additional hypotheses below are stated in terms of the original
arrows. Replacing them by existence of isomorphic grades would be weaker
and would not suffice to identify the crossing maps.
Assume cartesian lifts $c_{r,Q}:r^*Q\to Q$ over $r\in\J$, whose
universal property tests all total grade arrows. Assume a natural
transformation $\nu:d\Rightarrow Tp$, natural on all grade arrows, with
\begin{equation}\label{eq:nulaws}
 \nu_{U_B}\eta^L_B=\eta_B,
 \qquad \nu_{Q\di R}\mu^L_{Q,R}
                   =\mu_B(T\nu_Q)\nu_R.
\end{equation}
Require the following three tests:
\begin{align}
 \Delta_{r,Q}=(dc_{r,Q},\nu_{r^*Q}):d(r^*Q)
       &\iso dQ\times_{TB}TE;\label{eq:display}\\
 U&\text{ preserves cartesian display arrows, and}\notag\\
 (c_{r,Q},c_{dc_{r,Q},R})
       &\text{ accumulates to a cartesian arrow};\label{eq:cartalg}\\
 v_!r_u^*Q&\iso r^*u_!Q
       \quad\text{for every pure pullback of a display}.\label{eq:gradebc}
\end{align}
The last map is the canonical cartesian--opcartesian mate. The display
isomorphism implies $dc_{r,Q}\in\J$. Equation~\eqref{eq:gradebc} does
not restrict $u$ or $v$ to displays.

\paragraph{The two transport tests are independent.}
Even for finite surjections, with identity unit and multiplication bodies
and display automatic, neither \eqref{eq:cartalg} nor \eqref{eq:gradebc} implies the
other. For union grades on subsets of size at most three, cartesian
accumulation holds but a base-change lift would have to cover four
pairs across two fibres, and the mate fails. For grades of lists with
a duplication counter, base change holds but multiplying restricted
grades forces a duplicated letter to agree at both positions; restricting
the product permits different letters, so cartesian accumulation fails.
The companion gives both constructions and checks the one-sided laws
\cite{TransportIndependence}.

\begin{lemma}\label{lem:comma}
These hypotheses provide the relative content corners of
Definition~\ref{def:corners}, with occurrence $a_{(Q,f)}=\nu_Q f$.
\end{lemma}
\begin{proof}
Lift $(Q,f)$ along $r$ using grade $r^*Q$ and body the second projection
of $A_r=A\times_{dQ}d(r^*Q)$. A testing grade square factors uniquely
through $c_{r,Q}$, even when its remaining edge is not a display.
The source pullback then forces its left edge and body. This proves
cartesianity of the actual content square. Equation~\eqref{eq:display}
identifies its source with $A\times_{TB}TE$. Naturality of $\nu$ and
\eqref{eq:nulaws} prove all of \eqref{eq:arity}.

For a pure $u:B\to B'$ the lifted grade is $r^*U_{B'}\cong U_E$;
its source is $B\times_{TB'}TE\cong B\times_{B'}E$ by cartesianity
of $\eta$. Under the grade comparison the chosen cartesian arrow is
the original $U(r):U_E\to U_{B'}$. Naturality of the unit body gives
\[
 dU(r)\,\eta^L_Ev=\eta^L_{B'}rv=\eta^L_{B'}u r_u .
\]
The other display coordinate is $\eta_Ev$ by the occurrence unit law.
These are the two display coordinates of the lifted body, so it is
$\eta^L_Ev$ and the pure pullback square is canonical.

For horizontal stability let $(Q,f)$ be followed by $(R,b)$ and lift the
last arrow along $r:E\to C$. Set $R'=r^*R$, $j=dc_{r,R}$ and
$B'=B\times_{dR}dR'$ with projections $v:B'\to B$, $b':B'\to dR'$.
The grade of the horizontal paste is identified over its original
projection by
\[
 R'\di b'_!v^*Q\cong R'\di j^*b_!Q
                    \cong r^*(R\di b_!Q).
\]
The first map is \eqref{eq:gradebc}; the second is the canonical factor
from \eqref{eq:cartalg}. Their defining equations identify the pasted
grade arrow itself as cartesian. For $a=\nu_Q f$, $c=\nu_Rb$, the
source pullbacks compare by
\[
 A\times_{TB}TB'\cong A\times_{TTC}TTE
                          \cong A\times_{TC}TE,
\]
using preservation of pullbacks by $T$ and then cartesianity of $\mu$
at $r$. The second occurrence coordinate is flattening; the first display
coordinate is the actual composite body, because the pasted square is a
content square. This identifies the top content arrow with the direct
lift, not just its occurrence image. All intervening left edges are in
$\J$ by \eqref{eq:display}. Thus horizontal stability holds.
\end{proof}

\begin{corollary}[Recovery on the original maps]\label{cor:set}
For the finite polynomial resource calculus of Section~\ref{sec:resources}
and comma content satisfying \eqref{eq:nulaws}--\eqref{eq:gradebc},
Theorem~\ref{thm:glue} defines composition on
\[
                       H(A,C)=(e_A\downarrow d_C).
\]
It is the independent Set interaction of \cite{O1}, including its native
arrows, all pure balances, four structural comparisons and multiplication.
\end{corollary}
\begin{proof}
For comma content, \eqref{eq:normal} identifies $\grade_N$ with precisely
these comma categories and their pairs $(\alpha,\beta)$; no quotient of
native arrows is taken. To identify the crossing completely, let $P$ lie
over $\pi:E\to B$ and $Q\in\Lcat_B$. Put
\[
 \widehat P=L_{\nu_Q}P,\qquad \widehat Q=\pi^*Q,
 \qquad \chi=\Delta_{\pi,Q}^{-1}.
\]
At $y\in dQ$ with $\nu_Q(y)=(s,(b_i))$, the context $\widehat P$
has the independent choices and row tensor of \eqref{eq:occaug}.
Its carrier is $dQ\times_{TB}TE$ and $\chi$ is its body into
$d(\pi^*Q)$. Thus for $(Q,f)$ followed by $(P,g)$ the crossing is
\begin{equation}\label{eq:setcut}
 \left(f^*\widehat P,\ g_!\pi^*Q,\
 e(f^*\widehat P)\longrightarrow e\widehat P
    \xrightarrow{\chi}d(\pi^*Q)
    \xrightarrow{d\kappa_{g,\pi^*Q}}d(g_!\pi^*Q)\right).
\end{equation}
All of $\pi^*Q$ is retained, whether or not $f$ reaches every part of it.

For a grade arrow $m:Q\to Q'$ over $u$ and a decorated square with
carrier $h$, the grade map is the opcartesian factor
\begin{equation}\label{eq:nativegrade}
 h_!\pi^*Q\longrightarrow\pi'^*Q',\qquad
 c_{\pi',Q'}\widetilde m=m c_{\pi,Q}.
\end{equation}
Its companion context map is exactly \eqref{eq:occmap}, not a map chosen
from $h$ alone. The body equation follows by applying the two coordinates
of \eqref{eq:display}: they are $(dm,Th)$ applied to the old coordinates.
For a pure middle $u$, the same factorization identifies the grade
comparison with \eqref{eq:gradebc}; the body coordinates are
$(d\kappa_{u,Q}y,Tv(t))$ on a choice $(y,t)$. These identify the
product arrows and every generating balance of the independent law.

The four rows of \eqref{eq:four} also agree, with their original marks.
The content-unit row restricts the one-position comparison along
$\eta^L_B$ and uses $U_E\cong\pi^*U_B$. The context-unit row uses
forced singleton choices and $(\id_B)^*Q\cong Q$. The content-multiplication
row restricts leaf regrouping along $\mu^L_{Q,R}$ and uses the comparison
between $\pi^*Q\di(dc_{\pi,Q})^*R$ and $\pi^*(Q\di R)$. Its first
coordinate is $\mu^L_{Q,R}(z)$, not $z$. The context-multiplication row
uses the shuffle natural on \eqref{eq:dependent} and the comparison
$(\pi\rho)^*Q\cong\rho^*\pi^*Q$. Each grade comparison is the canonical
factor over its original projection. This is the complete component
identification, including the ten pastings already proved for those
components. Finally, for successive $(P,Q,f)$ and $(D,S,g)$ the new
context and grade are
\[
                 P\st f^*\widehat D,
                 \qquad S\di g_!\pi_D^*Q.
\]
The body is \eqref{eq:setcut} followed by the given content multiplication,
so the original multiplication is identified as well.
\end{proof}

\subsubsection{Receipts satisfy the tests}\label{sec:receipts}
The receipt calculus behind the two-resource merge (Section~\ref{sec:paired})
extends to any extensive base with pullbacks and a cartesian monad. Take
\[
 \Lcat_B=\B/(TB),\qquad dQ=B+Q,\qquad
 \nu_Q=[\eta_B,q:Q\to TB].
\]
Transport along $u$ postcomposes the label by $Tu$ and keeps $Q$;
restriction along $r$ pulls the tag object back along $Tr$.
The unit has tag object $0$. For $t:R\to T(B+Q)$ the accumulated grade
has tag object $Q+R$ and label
\begin{equation}\label{eq:receipts}
                      [q,\mu_B(T\nu_Q)t]:Q+R\longrightarrow TB.
\end{equation}
Its body is coproduct reassociation. Grade maps act on their actual tag
summands. For three stages the early tags are unchanged; the last-stage
labels agree by associativity and naturality of $\mu$. Coproduct
coherence proves the one-sided pentagon and units on tag maps, and
relabelling commutes with these operations. This is a valid content
calculus before any resource rig has been chosen.

Extensivity and cartesianity of $\eta$ give
\[
 E+(Q\times_{TB}TE)\cong(B+Q)\times_{TB}TE,
\]
which is precisely \eqref{eq:display} and its unit instance.
For accumulation, the old $Q$ summand pulls back directly. For the $R$
summand, apply $T$ to this pullback and then use the cartesian square of
$\mu$ at $r$. The result identifies restricted nested labels with the
pullback of their flattened labels, over the \emph{same original} $R$.
Consequently the accumulation comparison preserves its original projection
and is cartesian. The pure mate is invertible because $T$ preserves the
base pullback and transport keeps every tag. This verifies all three
tests and justifies every receipt and grade map used in the receipt
composite. It also gives the content input for the following indexed
example without altering the formulas.

\section{A change of stage that erases colours and merges witnesses}\label{sec:indexed}
Theorem~\ref{thm:glue} also applies when interfaces have two stages.
A transition between them can erase occurrence colours and merge several
rows into one. The resource lift must act on that noninvertible
transition, including its coefficients; the receipt construction of
Section~\ref{sec:receipts} supplies the content corners on the same base.
Let $\B$ be presheaves on $0\to1$, written
$X_1\to X_0$. Every natural transformation is pure; displays are the
pointwise finite-fibre maps. Fix a nontrivial group $G$, for example the
two-element group, and define
\begin{equation}\label{eq:coloured}
 (TX)_1=\coprod_{n\geq0}(G\times X_1)^n,
 \qquad (TX)_0=\coprod_{n\geq0}X_0^n.
\end{equation}
Restriction forgets the colours and restricts each entry. At stage 1 the
unit uses colour $1_G$, and multiplication concatenates words while
multiplying an outer colour with each inner colour in that order.
At stage 0 these are the ordinary list operations. Group associativity
and list concatenation prove the monad laws; erasure commutes with them.

This is a polynomial internal to the presheaf category. Its shape object
is $\coprod_nG^n\to\mathbb N$, with restriction recording length.
Positions are a shape together with an index $1\leq i\leq n$; restriction
keeps the index. Over a generalized shape at stage $c$, its position
object is a coproduct of $n$ copies of the representable $yc$, so its
labellings give precisely the stagewise formula \eqref{eq:coloured}.
The polynomial is not a constant Set polynomial: its length-one shapes
restrict by $G\to1$. Its unit and multiplication have bijective addressed
position maps at each stage. Their naturality squares are therefore
pullbacks, and $T$ preserves pullbacks stagewise. Finite products of
finite fibres give $T\J\subseteq\J$.

Use the same symmetric rig at both stages. A decorated display now
includes a stage-1 and a stage-0 decorated display and a structural
square between them of the form \eqref{eq:aggregation}. That square may
merge rows and have noninvertible coefficient maps. A map of such objects
must commute on those coefficient maps, not just on the presheaf carriers.
Horizontal composition uses \eqref{eq:dependent} stagewise.
For restriction along a natural transformation, take the two stagewise
cartesian lifts and induce their connecting square by the lower lift's
universal property. A test square factors at each stage; restoring the
lower cartesian arrow proves naturality of the factors. Thus the square
is cartesian in the category of diagrams, and these cartesian squares
remain stable under horizontal pasting.

Lift $T$ by coloured occurrence products at stage 1 and uncoloured
products at stage 0. On the structural restriction use
\eqref{eq:occmap} in the same positions, with the given aggregation
maps, while forgetting their colours. The composite-fibre calculation
for that formula proves functoriality on diagrams. Naturality of the
shuffle on \eqref{eq:dependent} makes the horizontal comparison a map
of diagrams. Erasure preserves the leaf bijections, so singleton and
flattening squares likewise commute with the structural restriction.
Their cartesianity and preservation of cartesian restrictions follow
from the stagewise factorization just proved. Only \emph{after} these
naturality checks are the remaining equations detected by stage
evaluation. This verifies Definition~\ref{def:lift} internally.
The receipts of Section~\ref{sec:receipts} provide content: all their
pullbacks, sums and monad maps are presheaf maps, so the proof of the
three tests applies without modification.

\subsection{A cut whose restriction is visible}
Let $B$ be terminal, and let the context carrier restrict by
$E_1=\{1,2\}\to E_0=\{*\}$. Upstairs give both rows weight $k=\mathbb Q$;
downstairs give the single row weight $W=k^2$. Both total augmentations
are the sum of coordinates. The structural aggregation $k\oplus k\to W$
can be either $1_W$ or $p(a,b)=(a+b,0)$. Both define lawful internal
contexts, but they are not isomorphic: an isomorphism of diagrams would
preserve the rank of the structural map.

Choose one receipt at each stage, restricting to the other, with upstairs
label $[(g,*),(1_G,*)]$ for $g\ne1_G$ and downstairs label $[*,*]$.
Use the generic content body $\id_{dQ}$ and identity context tail. At
that receipt the lifted grade has four upstairs tags, one for each
$(i,j)\in\{1,2\}^2$, and one downstairs tag. Its restriction is the
actual fold $4\to1$. The crossed context has four upstairs rows of
weight $k$ and one downstairs row of weight $W\otimes W$.
Formula~\eqref{eq:occmap} makes its structural coefficient map
\[
 k^4\cong W\otimes W
       \xrightarrow{\ 1\ \text{or}\ p\otimes p\ }W\otimes W.
\]
The ranks are four and one. In particular the $(2,2)$ input goes to
$e_2\otimes e_2$ in one case and $e_1\otimes e_1$ in the other.
The grade restriction and colour erasure are identical. Thus even the
same lists of stage objects do not determine the internal crossed object.
There is no isomorphism of presheaf contexts identifying these two
restrictions, even allowing arbitrary invertible changes of coefficient
coordinates. The calculation evaluates the actual cross-stage restriction
on coefficient maps and receipt tags.

\subsection{Changing the occurrence parameter}
Erasing the occurrence colours while leaving the two stages in place
defines a noninvertible change of the input calculus. The computation
in Appendix~\ref{app:colour} distinguishes this change of parameter
from both the stage restriction above and a pure change of interface.

\subsection{Why this does not permit arbitrary position restriction}
Polynomial positions and context witnesses have different jobs. In this
example the positions of a fixed shape restrict bijectively; its witness
rows need not. Pointwise finiteness alone cannot replace that position
condition. In the discrete rig $(\mathbb N,+,\cdot)$, let a single row
of weight $1$ map to the first of two rows of weights $(1,0)$. It is a
valid augmented additive transfer: its two fibre conditions are $1=1$
and $0=0$. But tensoring all witness rows would demand a map $1\to0$
after restriction. No such arrow exists in the discrete rig.
The construction avoids this obstruction by multiplying over polynomial
positions and using directed transfer for witnesses separately in each
position. An indexed norm doctrine with additional structure is not ruled
out; it is not supplied by finite fibres alone.

\section{What a crossing can force, and what it cannot}\label{sec:recognition}\label{app:recognition}
Theorem~\ref{thm:glue} constructs a law from independent inputs.
This section asks which tests a given crossing forces. The comparison
uses the original projection maps, so it starts with the comma
presentation of Section~\ref{sec:recovery} and marks those maps.

In the comma setting retain the one-sided calculi, chosen actual relative
grade lifts and \eqref{eq:nulaws}, but do not assume
\eqref{eq:display}--\eqref{eq:gradebc}. Call a crossing \emph{marked} here
if its generic output uses the complete carrier $dQ\times_{TB}TE$, the
grade $\pi^*Q$ with its original projection $c_{\pi,Q}$, and an invertible
body $\chi$ preserving both display coordinates. Its unit and accumulation
comparisons must preserve the original unit and accumulated
projections, together with the addressed context coordinates. Its
pure-middle comparison must preserve the original commuting
cartesian--opcartesian square and its factorization equations. These
requirements do not simply declare the comparisons canonical; they give
equalities from which that fact can be proved.

Let $\mathcal D\subseteq\J$ consist of projections realized, up to
isomorphism, by admitted contexts.
\begin{proposition}[Recognition with the original marks]\label{prop:recognition}
A marked crossing forces \eqref{eq:display} and \eqref{eq:cartalg} on
$\mathcal D$, and invertibility of the canonical mate
\eqref{eq:gradebc} for every pure pullback of a realized projection.
The marked grade comparisons are the canonical factors.
\end{proposition}
\begin{proof}
Preservation of the two coordinates gives $\Delta_{\pi,Q}\chi=1$.
As $\chi$ is invertible, $\Delta_{\pi,Q}=\chi^{-1}$, proving display.
The original unit arrow factors through the chosen cartesian lift;
so does the accumulation of the original cartesian pair. The projection
marks make the claimed structural comparisons factors of exactly those
same arrows. Cartesian uniqueness identifies each with the canonical
factor, whose invertibility proves cartesian preservation. The two
factorization equations for pure-middle balance similarly identify its
comparison with the canonical mate, so that mate is invertible.
\end{proof}
If the body is not invertible, the coordinate equation proves only that
$\Delta$ has a section. If a display cannot be realized, no generic
crossing tests it. This proposition therefore does not recognize an
arbitrary unmarked law, nor reconstruct the independently lifted monad
from one. It is a statement about preservation of specified evidence.

For the Set resource calculus, the realized class can be determined.
Assuming choice, a nonempty finite fibre is always realizable: put $I$
at one point, $0$ at the others, and augment using the additive unit maps.
An empty fibre is realizable exactly when $\V(0,I)\ne\varnothing$.
Consequently
\begin{equation}\label{eq:realized}
 \mathcal D=
 \begin{cases}
 \text{finite-fibre members of }\J,&\V(0,I)\ne\varnothing,\\
 \text{finite-fibre surjections in }\J,&\V(0,I)=\varnothing.
 \end{cases}
\end{equation}
Both classes are closed under pullback and composition. A fibre of $Tf$
is a finite product of fibres of $f$, including a singleton for the empty
product, so $T$ preserves these classes. Using $\mathcal D$ in the
sufficient theorem leaves every admitted context and every native map
between them unchanged. It restricts the required lifts, not the maps
against which their cartesianity is tested. Thus the same tests are
sufficient for that entire context calculus and necessary for its marked
crossing, with the stated scope.

\subsection{A relative construction need not extend globally}
There is a concrete obstruction to replacing relative lifts by global
co-crossedness. Define $Z(X)=\varnothing$ for $X=\varnothing$ and
$Z(X)=1$ otherwise. The endofunctors $\id,Z$ form a monoid under
composition and give a constant discrete grading, with extension by
evaluation. It is a one-sided content calculus. With $T=\mathsf{List}$,
assign the identity grade the singleton occurrence and the $Z$ grade the
empty occurrence. Unit and substitution hold, including at the empty set.
Grade lifts retain their grade; their mates and cartesian accumulation
comparisons are automatic.

Display holds along every surjection: its source and target have the
same inhabitation status, and the only list above an empty list is empty.
It fails for $\varnothing\to1$ on the $Z$ grade, where display would
require an isomorphism
\[
        \varnothing\longrightarrow
 Z(1)\times_{\mathsf{List}(1)}\mathsf{List}(\varnothing)=1.
\]
More directly, let $i_0,i_1:1\to2$ select different points. Compose the
$Z$ arrow $1\rightsquigarrow1$ with pure $i_0$. Lifting this composite
along $i_1$ gives source $1\times_{Z(2)}Z(1)=1$. Lifting $i_0$ first
along $i_1$ gives source $\varnothing$, and then lifting the $Z$ arrow
along $\varnothing\to1$ leaves source $\varnothing$. The horizontal
paste is not cartesian. This violates the global hypothesis of the
corners construction, but does not affect the law on finite-fibre
surjections. Both $i_0$ and $i_1$ are still pure maps in its tensor.

The realized-class distinction is also substantive. In the discrete rig of nonnegative real numbers, an augmentation
requires the weights to sum to $1$. There is no arrow $0\to1$, so only
nonempty fibres can be augmented. The relative construction above applies. In
finite-dimensional vector spaces with all linear maps, $0\to k$ exists,
so the same content fails the marked display test on a realizable empty
fibre. These are different context doctrines, not contradictory answers
to one unchanged problem. The obstruction applies to the marked normal
form, not to every possible representation of an interaction.

\subsection{The claimed boundary}
The output is a construction on native hom categories, not a free or
initial object among all assemblers. Pullback uniqueness, the local
category pullback \eqref{eq:pullcat}, and the pseudocoend each give a
specific universal property. None supplies a category of all mixed
realizations with an initial object. Such a statement would need its own
comparison morphisms and proof.

Likewise, the occurrence assignment in \eqref{eq:arity} is strict on
native cells: arbitrary arity-changing content maps are not included.
The theorem does not supply shared-choice semantics, feedback, infinite
tensor products, lax noninvertible distributors, arbitrary internal
position maps, or a higher-categorical bispan enhancement. The coloured
presheaf example establishes a particular internal generality with
noninvertible witness restriction and noninvertible parameter change;
it does not erase these restrictions. The Set recovery and marked
recognition have different conclusions and have been kept separate
for that reason.

\section{The local occurrence equations}\label{app:local}
The action $L_aP=a^*\Tt P$ in \eqref{eq:L} pulls a lifted resource display
back along an occurrence map. The following calculation proves its
comparison equations \eqref{eq:IU}--\eqref{eq:middle} using only the
decoration-side lift of Definition~\ref{def:lift}. For a decorated square over $(h,u)$ and
Kleisli arrows satisfying $Tu\,a=a'k$, apply $\Tt$ to the square and
precompose with the cartesian arrow $a^*\Tt P\to\Tt P$. Factoring
through $a'^*\Tt P'$ gives the action of $L$; its carrier coordinates
are $(k,Th)$. Cartesian uniqueness proves its identity and composition
laws on actual decorated maps.

For $K$ the construction is
\[
 a^*\Tt(b^*\Tt P)
 \cong a^*(Tb)^*\Tt^2P
 \cong (Tb\,a)^*\mu^*\Tt P
 \cong L_{b\bullet a}P.
\]
The first comparison uses preservation of cartesian squares and the
second the cartesian square of $\bar\mu_P$. Its projection into $TE$
is Kleisli flattening, proving the typing equation after
\eqref{eq:middle}. The analogous construction with $\bar\eta_P$ gives
$U$. For $F$, use the horizontal compositor of $\Tt$ and observe that
restriction of $\Tt P\st\Tt S$ along $a$ is the horizontal paste of
restriction of $\Tt P$ along $a$ and of $\Tt S$ along $a_P$.
Horizontal cartesian stability makes that paste the required lift.
The horizontal unit comparison similarly gives $I$. Each construction
is natural in every square in its input diagram.

There is a reliable equality test for their composites: restore the
cartesian arrows from each restriction to the original display. The two
routes in associativity of $K$ then end by
$\bar\mu_P\Tt\bar\mu_P$ and $\bar\mu_P\bar\mu_{\Tt P}$.
Naturality moves intervening restriction arrows to the common source,
and the monad equality identifies those endings. The two $K$ unit
identities similarly use the two monad unit equalities. The restored
$F$ associativity diagram is the horizontal functor pentagon of $\Tt$;
its units are the functor unit equations.

For the mixed equation \eqref{eq:middle}, before restriction along
$(Tb)a$ the restored square is
\[
 \begin{array}{ccc}
 \Tt^2(P\st S)&\longrightarrow&\Tt^2P\st\Tt^2S\\
 {\scriptstyle\bar\mu_{P\st S}}\downarrow&&
      \downarrow{\scriptstyle\bar\mu_P\st\bar\mu_S}\\
 \Tt(P\st S)&\longrightarrow&\Tt P\st\Tt S.
 \end{array}
\]
The top arrow is the composite double-functor comparison. This square
is exactly the horizontal-composition axiom for the vertical double
transformation $\bar\mu$. Its horizontal-unit variant uses the unit
axiom of that transformation. Replacing $\bar\mu$ by $\bar\eta$
gives compatibility of $F,I$ with a pure occurrence unit. Cancellation
through the final cartesian arrows proves the original equations in
$\D_1$. In particular no faithful functor from decorations to carriers
has been assumed.

\section{Descent and all ten pastings}\label{app:coherence}
The four comparisons in \eqref{eq:four} were described on
identity-tail representatives. They must respect the pure balances
that change the middle interface, and they must satisfy the ten
pseudodistributive pastings. A carrier bijection alone does not
identify the maps in either calculation.

An internal tail in a string of contexts can be moved by the invertible
balance
\[
 [(P,g),(S,f)]\cong[(P,\id),(g^*S,fv)],
\]
where $v$ is the pullback map to the carrier of $S$. A tail next to content
moves by pure precomposition; the last tail remains a pure output action.
These are isomorphisms in a category-valued pseudocoend. They reduce
objects to the representatives used in \eqref{eq:four}, not to equal
objects in a set quotient. Naturality on product arrows follows on
content from \eqref{eq:factor} applied to the pasted square, and on
contexts from naturality of the local occurrence comparisons.

The remaining balance checks have three seam types. At an $N|N$ seam,
inserting a pure $u$ puts a pure column between two content columns.
Cutting on either side gives the same three-column paste; the middle
cartesian square is the canonical pure pullback. Its factor is
\eqref{eq:mate}. Horizontal associativity followed by cartesian uniqueness
identifies the content routes. On contexts this is associativity of $K$
with middle Kleisli arrow $\eta u$, followed by
$U_u:L_{\eta u}P\cong u^*P$ for the display $P$ at that seam.

At an $N|M$ seam, the pure terminal column can instead restrict the first
display. The comparison is exactly the defining mate square. In a
vertical stack, the pullback of the composite display is the vertical
paste of its pullbacks; in a horizontal string, the pure column pastes
with the final content column. After restoration and unbending, both
routes are the same content square. On contexts these statements are
respectively compatibility of $U$ with $F$, and the $K$ composition law
for a terminal pure arrow, together with cartesian preservation by $\Tt$.

At an $M|M$ seam, moving the tail is the displayed slide to $g^*S$.
Interchange with its pure pullback square identifies the two large
content diagrams, and cartesian cancellation identifies the factors.
On contexts this is naturality of $F$ for that cartesian restriction,
using $L_{\eta g}\cong g^*$. Endpoint versions are the first and last
slides or output transport. Consecutive slides compose by mate pasting;
disjoint slides commute by functoriality. This proves descent of all four
comparisons to the relative tensors, including unrestricted forward tails.

Walker gives equivalent presentations and a smaller coherence test
\cite[Thm.~5.0.1 and Sec.~5.3]{walker}; the full list is retained here
to identify the original constraints. For completeness, the following
table uses C1--C8 of \cite[Def.~4.1]{GL} and the derived pastings
C9--C10 of \cite[App.~A]{GL}, with their $S=N$ and $T=M$.
Its equations are understood on the common representative, with the
canonical restriction and bracketing comparisons inserted. The table is
an index of proof obligations, not a substitute for the cancellation and
local arguments that discharge them.
\begin{center}\small
\begin{tabular}{@{}cp{.27\linewidth}p{.59\linewidth}@{}}
\toprule
&Restored content rectangle&Decorated equality\\\midrule
C1&Three columns, one row&Associativity of $K$, from the two monad multiplications on $\Tt^3P$.\\[3pt]
C2&Left unit column&Left unit of $K$, using one monad unit identity.\\[3pt]
C3&Two columns, identity row&Horizontal-unit axiom of $\bar\mu$, comparing flattening with $I$.\\[3pt]
C4&Both units&Horizontal-unit axiom of $\bar\eta$.\\[3pt]
C5&Two columns, two rows&Equation~\eqref{eq:middle}, the horizontal-composition axiom of $\bar\mu$.\\[3pt]
C6&Identity column, two rows&Horizontal-composition axiom of $\bar\eta$, identifying $F$ at a pure unit with ordinary restriction composition.\\[3pt]
C7&One column, three rows&Associativity of $F$, from the horizontal compositor equation of $\Tt$.\\[3pt]
C8&Final identity row&Right unit of $F$, with the corresponding horizontal unitor.\\[3pt]
C9&Right unit column&Right unit of $K$, using the other monad unit identity.\\[3pt]
C10&First identity row&Left unit of $F$, with the other horizontal unitor.\\\bottomrule
\end{tabular}
\end{center}
For C5, take two consecutive content columns $n,m$ and two display
rows $P,S$, the latter with carriers $E\to B$ and $F\to E$. One route compares the horizontal pastes before the vertical paste;
the other compares the vertical pastes before the horizontal paste. On content the comparisons restore to the
same outer cartesian square: the two-by-two paste over $F\to B$ with
bottom $nm$. Both routes retain the original source projection and the
final occurrence projection to $TF$, with the nested occurrence
coordinates flattened by $\mu$. These coordinates match the boundaries;
cartesian uniqueness in the content calculus identifies the cells.
On resources, restoring the restriction identifies the comparison with
the horizontal-composition square for $\bar\mu$ displayed in
Appendix~\ref{app:local}; the two routes are precisely
\eqref{eq:middle}. Equality here is on decorated maps, not a consequence
of equality of carrier coordinates.

For each content entry, fill the indicated rectangle from the lower
right. A comparison of a subrectangle, whiskered and then compared on the
whole rectangle, satisfies the same equation
$C\,t=\text{restored paste}$ as direct comparison of the whole rectangle.
Interchange and naturality of the horizontal associator justify this
equality before cancellation. Empty rows are vertical identities; empty
columns are horizontal identities on displays. When bracketings differ,
the bottom and right one-sided associators and unitors have the same
composite by their pentagon and triangles. Thus cartesian uniqueness
proves each content equality, including inverse orientations.
Appendix~\ref{app:local} proves each decorated equality.

The preceding seam checks transport these component equations to all
representatives. Rebracketing iterated pseudocoends represents the same
multiply pseudobalanced cocone, so its associators are precisely the
slides already treated. Interchangers for disjoint slots are naturality
squares, not additional coefficient equations. Equality of components
on all objects is equality of modifications. Finally equality of both
native components is equality in $\grade_N$, and the equivalence
\eqref{eq:normal} reflects it. This completes the ten-pasting proof.

\section{Changing the independent inputs}\label{app:parameters}
A change of resource or content model can be compared before or after
the crossing \eqref{eq:cross}. This appendix specifies the maps of
inputs for which those two routes agree up to the stated comparison.
A compatible input map from unprimed to primed data consists of three pieces.
First, a pullback-preserving $V:\B\to\B'$ takes $\J$ into $\J'$, and
$\theta:VT\to T'V$ is cartesian on displays and satisfies
\[
 \theta V\eta=\eta'V,\qquad
 \theta V\mu=\mu'V\,(T'\theta)(\theta T).
\]
Second, a pseudo-double functor $D:\D\to\D'$ over $V$ preserves
carriers and codomain-cartesian squares, and a vertical double
transformation $q:D\Tt\to\overline{T'}D$ over $\theta$ has cartesian
display squares and satisfies
\[
 qD\bar\eta=\bar\eta'D,\qquad
 qD\bar\mu=\bar\mu'D\,(\overline{T'}q)(q\Tt).
\]
Third, a strong pseudofunctor $G:N\to N'$ over $V$ on objects has a
compatible invertible pure icon $G\pure\cong\pure'V$. Its induced
functor on bent squares preserves the specified content cartesian squares,
and $a'_{Gn}=\theta_BVa_n$. These data compose by ordinary functor
composition and pasted transformations and icons. Their equations are
stable under this composition; cartesian squares paste. Fixing bracket
conventions gives a category of such inputs.

The induced context functor takes $(P,g)$ to $(DP,Vg)$. Its strength is
cartesian preservation followed by the horizontal compositor of $D$.
A relative tensor representative with middle $B$ goes to one with middle
$VB$, and its balance along $u$ goes to the balance along $Vu$.
The comparison between the two orders of crossing and changing inputs
has decorated component
\[
 D(a_n^*\Tt P)\cong(Va_n)^*(D\Tt P)
 \cong(Va_n)^*\theta_B^*(\overline{T'}DP)
 \cong L'_{a'_{Gn}}DP.
\]
Its content component compares $GC_r(n)$ with $C'_{Vr}(Gn)$ and then
transports along $Vg$. Their carrier maps agree: both are the specified
pullback comparison, with the occurrence projection transported by
$\theta_E$. This is where cartesianity of $\theta$ and $q$, rather than
invertibility of $\theta$ itself, is needed.

Restoring the content cartesian squares proves naturality on
\eqref{eq:factor}; the pure icon and opcartesian factorization prove it
on \eqref{eq:mate}. For the four comparisons the decorated equations
are respectively preservation of the occurrence unit, horizontal unit,
occurrence multiplication and horizontal compositor. On content the
restored squares are the images of the same unit, identity, horizontal
paste and vertical paste, hence their factors agree. For successive input
maps, the decorated comparison restricts the pasted $q$, while the
content comparison has the same restored projection as the paste of the
two comparisons. This proves compatibility with composition and identity,
and therefore the asserted strong map of composite calculi. No mixed-law
map was assumed as input. The colour-erasure map of Section~\ref{sec:indexed}
is a noninvertible instance of exactly this specification.

\subsection{A computed change of parameter}\label{app:colour}
A change of occurrence model can forget a job's colour while preserving
its address. Work over two-stage presheaves \(X_1\to X_0\), with a
single symmetric rig of coefficients at both stages. The monad \(T\)
forms coloured lists at stage~1 and ordinary lists at stage~0, as in
\eqref{eq:coloured}. Let \(T'\) form ordinary lists at both stages.
Define \(\theta:T\to T'\) by forgetting colours at stage~1 and acting
identically at stage~0. Erasure commutes with unit and concatenation,
including multiplication of nested colours before erasure, so $\theta$
is a monad map. Its naturality square at a carrier map is a pullback:
given the coloured image word and an uncoloured lift, give each entry
of that lift the colour of the corresponding image position. This
constructs the unique coloured lift. The comparison on decorated
displays retains the row objects and uses the identity row maps under
this fibre bijection; hence it preserves the shuffle, unit and
flattening. On content the comparison postcomposes receipt labels
with $\theta$, keeping their tag maps and bodies. The same pullback
argument identifies the relative lifts over their original projections.
Thus this is an input map of Appendix~\ref{app:parameters}, with
identity base and decorated-display functors.

For the generic body of Section~\ref{sec:indexed}, the upstairs cut at
the receipt has four elements
\[
 (q,[(g,i),(1_G,j)])\longmapsto(q,[i,j]),
 \qquad i,j\in\{1,2\}.
\]
The formula is at once the carrier comparison and the tag comparison
over the receipt. It is bijective on this fibre and acts by the identity
on each coefficient $k\otimes k$. Downstairs the comparison is the
identity. Thus the comparison is invertible for this cut even though
the input change is not: the words $[(g,*)]$ and $[(1_G,*)]$ have the
same image.

To see naturality on noninvertible native maps, take the source
display with upstairs two $k$-rows and downstairs one $W=k^2$-row,
with the structural coefficient identification $k\oplus k\cong W$.
Fold it onto a display with one $W$-row at both stages and identity
restriction. Take the aggregation upstairs and downstairs respectively
to be $(\gamma,\gamma)$, under the displayed identification, with
$\gamma=1_W$ or $\gamma=p$ from Section~\ref{sec:indexed}. These are
maps of diagrams: the source structural coefficient identification
intertwines the two copies of $\gamma$, and both preserve the sum
augmentation. In either order of applying the crossing and forgetting
colours, the cut sends the four upstairs choices to one choice,
the four upstairs tags $q_{ij}$ to one tag, and the sole downstairs
tag to itself. The coefficient map from the four upstairs rows is
\[
 k^4\cong W\otimes W\xrightarrow{\ \gamma\otimes\gamma\ }W\otimes W .
\]
It sends the $(2,2)$ input to $e_2\otimes e_2$ for $\gamma=1_W$ and
to $e_1\otimes e_1$ for $\gamma=p$. The downstairs map is the same
tensor $\gamma\otimes\gamma$ on its sole row. Reindexing along $\theta$
changes neither calculation; it removes colours, not addresses.
The stage restriction is a different map, with four source tags
upstairs and one downstairs. A pure middle balance is different again:
it moves a map of interfaces inside a single calculus. This witness
compares two calculi on the same base, with no pure-middle map inserted.

\section{Adaptive destructive diagnostics}\label{app:diagnostic}\label{sec:diagnostic}
A diagnostic can consume a specimen and still have to decide which
assays should use the survivors. The worked example below uses
finite planar trees with terminal content grades and tracks
probability and cost on the resource side
\cite[Secs.~2--3]{Casebook}. Corollary~\ref{cor:set} supplies the
composition of tests with their result-dependent follow-ups.
The replacement claim in this example uses a separate coupling and
cost calculation.

Three specimens $r_0,r_1,r_2$ initially share one hidden fair quality
bit $q$. The screen consumes $r_0$. A low result gives a single
continuation holding both surviving specimens; a high result gives two
ordered continuations, one for assay $A$ on $r_1$ and one for assay
$B$ on $r_2$. On the low branch the confirmation $C$ consumes $r_1$.
A pass or fail sends $r_2$ to disposal; an uncertain result sends it
to rescue assay $D$. The continuation at each hole receives its
actual specimens. For example, the two high continuations do not
receive two copies of the original bundle.

\begin{center}\small
\begin{tabular}{@{}lcc@{}}
\toprule
Response law & $q=0$ & $q=1$\\
\midrule
$\Pr[\text{screen is high}\mid q]$&$1/4$&$3/4$\\
$\Pr[A=1\mid q]$&$1/4$&$3/4$\\
$\Pr[B=1\mid q]$&$1/8$&$7/8$\\
$C$ (pass, fail, uncertain)&$(6,1,3)/10$&$(1,6,3)/10$\\
$\Pr[D=1\mid q]$&$1/5$&$4/5$\\
\bottomrule
\end{tabular}
\end{center}
At a high result the resource calculation supplies one child witness
at each actual hole. The weights multiply \emph{conditional on the full
private prefix}, and the content calculation grafts both outcomes onto
the screen node. Averaging happens after this product. In particular,
writing $S_1$ for the high result, Bayes' formula gives
$\Pr[q=1\mid S_1]=3/4$, and hence
\[
 \Pr[A=B=1\mid S_1]
 =\tfrac14\,\tfrac14\,\tfrac18+\tfrac34\,\tfrac34\,\tfrac78
 =\tfrac12.
\]
If a rule hides the shared bit before filling the two holes and
multiplies their posterior marginals, it gives
$(5/8)(11/16)=55/128$ instead, a deficit of $9/128$ on this public
observation. Grafting retains the two labelled outcomes; the retained
private input supplies the dependence between them. The calculation
does not give the running program access to $q$ or a conditioning
oracle.

The casebook compares two implementations of the same controller.
Each nondisposal method in the reference implementation adds one fair
private audit bit. At an outcome of candidate weight $p$, the two
reference witnesses have weight $p/2$ each. Forgetting the audit bit
is a native map: summing its fibre recovers $p$, its bodies agree
after projecting away private audit state, and its cost can decrease.
For the high branch, for example, both child costs contribute:
\[
 (2,1)+(3,1)+(4,2)=(9,4)
 \quad\text{versus}\quad
 (4,2)+(5,2)+(7,3)=(16,7).
\]
The low-conclusive and low-rescue masses are $7/20$ and $3/20$,
and the high mass is $1/2$. Their candidate cost pairs are $(7,3)$,
$(13,5)$, and $(9,4)$, versus $(12,5)$, $(21,8)$, and $(16,7)$
for the reference. Consequently the expected costs are
$(89/10,19/5)$ and $(307/20,129/20)$. These are comparisons of
whole-run costs; they do not assert equality of exact-cost laws.

The composition theorem supplies the addressed grafting and its action
on the original maps. The casebook's separate comparison couples the
implementations on related full private inputs, preserves public causal
trees and actual returned capabilities, and bounds candidate cost
coordinatewise for its admitted finite closing clients
\cite[Sec.~3]{Casebook}. Here the three-specimen calculation checks
one controller and one shared-quality initialization. The broader
comparison needs its own assumptions on initialization, ownership,
local kernels and client observations.

\section{Robot transactions: what undo can preserve}\label{app:robot}

A robot that turns a fixture and then turns it back can still undo
someone else's work. Suppose the fixture starts at orientation~\(0\).
The robot turns it a quarter turn, an outside command resets it to~\(0\),
and the robot performs its inverse turn. The fixture ends at~\(3\).
The outside reset alone would have left it at~\(0\). A different failure
is possible even when the final state is right: a camera that saw the
quarter turn still has its record.

The finite model below makes these two failures explicit. The
commands have signatures that determine which continuations are
available after each outcome. Corollary~\ref{cor:set}, specialized to
the finite probability-and-cost trees in Section~\ref{sec:observations},
composes the commands on their resource and cost data. The primary
semantics proof is linked in Appendix~\ref{app:reuse}. A separate
interference argument establishes when an abort
preserves outside changes and when the whole execution looks atomic
to the admitted observers \cite{RobotTransactions}. Compensating
transactions can satisfy many application-specific recovery criteria
\cite{KorthLevySilberschatz}. Here the reference preserves the actual
outside actions and specifies which observations are compared.

\paragraph{Which commands can follow an outcome?}
The fixture state is \((q,h)\in\mathbb Z/4\mathbb Z\times\mathbb Z/3\mathbb Z\):
orientation \(q\) and an outside counter \(h\). Acquiring the fixture
returns either \(\mathsf{busy}\), with no journal, or a held lock and
a single owned journal, initially empty. The journal records commands
rather than the pose. A successful rotation \(a(q,h)=(q+1,h)\) appends
\(a\) to the journal; a successful undo applies \(a^{-1}\) and removes
that entry. Abort requires the empty journal. Commit and abort release
the lock and return journal authority to the manager. If a step
faults, the manager retains the lock in quarantine and reports an
unrecovered fault. Successful commit and abort continuations are
unavailable; the physical fault update requires its own specification. Successful undo
presupposes that the inverse can actually run; the equality below
covers successful reversible executions.

These commands use the finite addressed-tree instance of
Corollary~\ref{cor:set}, with terminal content grades and probability
and cost on the context side. At each leaf of a command's result tree,
the continuation receives that leaf's observation and capabilities.
In particular, the busy leaf cannot supply a journal to an undo.
In a composite, a prefix witness of mass \(p\) and cost \(c\) is paired
with one continuation witness at each addressed leaf \(b\). The
resulting mass and cost are \(p\prod_b p_b\) and \(c+\sum_b c_b\).
The body grafts each chosen continuation at its own address. A finite
catalogue and a bound on the number of steps keep the construction
finite. In the fixture, each rotation or inverse rotation has unit
cost, and deterministic operations have mass one.

\paragraph{When does compensation preserve outside changes?}
Each robot operation \(a_k\) is a permutation, including the undo
steps. Write \(P_k=a_k\cdots a_1\) for the prefix; composition is read
from right to left. The outside actions \(e_k\) need only be functions.
For a fixed history the reference applies the same labelled outside
functions to the untouched initial state. When \(e_k\) follows step
\(k\), final compensation by \(P_m^{-1}\) yields that reference state
for every initial state precisely when
\begin{equation}\label{eq:robot-full-prefix}
 (P_m^{-1}e_mP_m)\cdots(P_1^{-1}e_1P_1)=e_m\cdots e_1 .
\end{equation}
To see this, move each outside action past the prefix that precedes
it. In a completed abort \(P_m=\id\), so the final compensation in this
formula is the identity. If the admitted actions at each position
include identity and can be chosen independently, the equality holds
for all histories exactly when each allowed \(e_k\) commutes with
\(P_k\). Commutation makes each conjugate equal to \(e_k\); taking
identities at all other positions proves necessity. This condition
includes interruptions during undo.

The translations \(e_{r,s}(q,h)=(q+r,h+s)\) commute with the robot's
rotations, including translations that change the orientation.
The word \(a,a,a^{-1},a^{-1}\) therefore ends at the state obtained by
the same outside translations alone. The reset in the opening
paragraph fails this test: its output after a rotation does not
carry the corresponding rotation of its output before the rotation.

\paragraph{What can a client observe?}
While the lock is held, allow the client and scheduler to observe
\(h\), the outside action labels and public controller data. The
abstract transaction records the same acquisition and release
responses, and delays physical motion until release. It applies
\(P_m\) on commit and identity on abort. Compare the two executions
with the same random coins and with decisions made from these common
observations. At position \(k\) before release, the concrete state is \(P_ky_k\) and
the abstract state is \(y_k=e_k\cdots e_1(x)\). Rotation leaves \(h\)
unchanged, so the two executions can make the same next decision.
At release the full physical state, ownership and response agree.
The comparison retains acquisition and release responses, the permitted
observations, outside action labels, and the full state at release. It
erases the robot's micro-events and elapsed time. The coupled executions
therefore have the same projected trace law. Equality extends through
post-release continuations that use the same admitted method kernels. For a general fixture the corresponding
hypotheses are \(oP_k=o\) for the permitted observation map \(o\), and
commutation for the allowed outside actions at each visited prefix.

A camera that can see orientation would distinguish the executions
immediately. A scheduler that sees pose can also choose different
outside translations in the two runs, even though each translation
commutes with the rotation. That is why the comparison specifies both
the primitive actions and the information used to choose them.
The native tree keeps the actuator events that the observation
projection erases. An abort with \(n\) turns and \(n\) inverses still
records \(2n\) steps. A deployment would have to establish the stated
motion, observation and authority hypotheses.

\section{Sources, reconstruction and verification}\label{app:reuse}
The public
\href{https://github.com/eternal-november/eos-research-programmes/tree/main/projects/contextual-composition}{contextual-composition project}
contains the underlying constructions, application proofs and formalization.
Each \textup{[PDF, TeX]} pair gives both formats at the same repository
revision. The Lean links identify a separate pinned formalization snapshot.

\paragraph{Recovering the construction.}
For the common algebra, begin with \cite{U}, in particular
\texttt{thm:coend}, \texttt{thm:mixed}, \texttt{thm:problaw} and
\texttt{thm:bicyclic}. They distinguish normalization, composition from a
supplied law, the concrete conservative example, and the obstruction to
universal conservative assembly. For constructing the law, use
\cite{Gluing} and \cite{O1}; Definitions~\ref{def:lift} and
\ref{def:corners} in this paper reproduce the relevant input contracts,
and Appendix~\ref{app:coherence} preserves the complete pasting check.
The marked converse uses only the realized display class
\eqref{eq:realized}. The previous comparative analysis is \cite{Comparison}.
For applications use the casebook \cite{Casebook} and the separate robot
source \cite{RobotTransactions}; their observation, ownership and
interference premises remain additional to the crossing theorem.

The finite-semantics proofs appear as \texttt{thm:runs} and \texttt{thm:comparison} in
\emph{Plural Semantics for Result-Dependent Protocols} \mbox{[\href{https://github.com/eternal-november/eos-research-programmes/blob/10e771644cfdf58c3d05e989e4b26cae7164c80a/projects/contextual-composition/protocols/plural-semantics/Plural_Semantics_for_Result_Dependent_Protocols.pdf}{PDF}, \href{https://github.com/eternal-november/eos-research-programmes/blob/10e771644cfdf58c3d05e989e4b26cae7164c80a/projects/contextual-composition/protocols/plural-semantics/Plural_Semantics_for_Result_Dependent_Protocols.tex}{TeX}]}.
The all-context comparison is
\emph{Contextual Refinement and Controlled Recovery}~\mbox{[\href{https://github.com/eternal-november/eos-research-programmes/blob/10e771644cfdf58c3d05e989e4b26cae7164c80a/projects/contextual-composition/protocols/refinement-recovery/Contextual_Refinement_and_Controlled_Recovery.pdf}{PDF}, \href{https://github.com/eternal-november/eos-research-programmes/blob/10e771644cfdf58c3d05e989e4b26cae7164c80a/projects/contextual-composition/protocols/refinement-recovery/Contextual_Refinement_and_Controlled_Recovery.tex}{TeX}]},
\texttt{thm:replacement}, with the primitive, pure-syntax and code-library
premises stated in its local input contract. The casebook explains their
application; it is not a substitute proof of those general finite results.

The
\href{https://github.com/eternal-november/eos-research-programmes/blob/main/projects/contextual-composition/NAVIGATION.md}{project navigation}
organizes these different reading tasks. The
terminology note~\mbox{[\href{https://github.com/eternal-november/eos-research-programmes/blob/10e771644cfdf58c3d05e989e4b26cae7164c80a/projects/contextual-composition/foundations/terminology/Interactad_Terminology_Note.pdf}{PDF}, \href{https://github.com/eternal-november/eos-research-programmes/blob/10e771644cfdf58c3d05e989e4b26cae7164c80a/projects/contextual-composition/foundations/terminology/Interactad_Terminology_Note.tex}{TeX}]}
records the presented-package usage of interactad. These are source-reading
routes, not substitutes for the proofs. A change of base, display class,
grade-map convention, occurrence functor or observation interface requires
renewed checking of the corresponding comparison, rather than inheritance
from a nearby example.

\paragraph{Formal verification.}
The
\href{https://github.com/eternal-november/eos-research-programmes/blob/fc2d2648058191b79b297a15161e85605cb13bb8/projects/contextual-composition/foundations/relative-construction/lean/README.md}{Lean package README}
and its
\href{https://github.com/eternal-november/eos-research-programmes/blob/fc2d2648058191b79b297a15161e85605cb13bb8/projects/contextual-composition/foundations/relative-construction/lean/SOURCE_CROSSWALK.md}{source crosswalk}
identify the actual declarations and their hypotheses at this revision.
The Lean~4.35.0 package includes fixed monad/comonad recoveries, Type-based
reader/writer and parameterized instances, and a mixed-category construction
from a supplied law. It also contains substantive general-grade
normalization results: primitive-balance localization with an exact
native-map inversion boundary, and selected endpoint and target
transport/coherence comparisons. These are stronger than object-level
or terminal-grade checks. They are still not a certificate for the entire
global relative factory, its iterated mixed tensor and full modifications,
or the general independent-input gluing theorem.

The package records scoped replay and axiom-audit evidence. This manuscript
revision did not rerun that Lean package; the verification claims here are
identified by the linked declarations and their existing evidence.
At the time of writing, the paper's general normalization and gluing
theorems have written proofs. Lean covers the subresults listed above;
the example checks verify their advertised finite calculations. The project will
formalize the remaining global normalization and gluing results and link
the verified development from the repository. This includes the global
endpoint packaging, native crossing maps, arbitrary pure balance and all
structural modifications.

\paragraph{Reproducing the finite examples.}
The \href{https://github.com/eternal-november/eos-research-programmes/blob/cdf5077225ded75a4f8042909a48ffd04d7cb51d/projects/contextual-composition/interaction/independent-criterion/check_gluing_witnesses.py}{Python checker} exercises colour erasure, dependent flattening, and selected
rational coefficient maps. Run it with Python 3 from the corresponding
checkout:
\begin{verbatim}
cd projects/contextual-composition/interaction/independent-criterion
python3 check_gluing_witnesses.py
\end{verbatim}
Appendix~\ref{app:coherence}
gives the general calculation; the named
\hyperref[par:dependent-fubini-maps]{Dependent Fubini on maps} paragraph
records the coefficient-map comparison.

\paragraph{Reading the later extensions.}
The
barrier feedback assembly route~\mbox{[\href{https://github.com/eternal-november/eos-research-programmes/blob/10e771644cfdf58c3d05e989e4b26cae7164c80a/projects/contextual-composition/foundations/native-whole-scope-feedback/Barrier_Feedback_Assembly_Reading_Route.pdf}{PDF}, \href{https://github.com/eternal-november/eos-research-programmes/blob/10e771644cfdf58c3d05e989e4b26cae7164c80a/projects/contextual-composition/foundations/native-whole-scope-feedback/Barrier_Feedback_Assembly_Reading_Route.tex}{TeX}]}
and
nonblocking recurrent assembly route~\mbox{[\href{https://github.com/eternal-november/eos-research-programmes/blob/10e771644cfdf58c3d05e989e4b26cae7164c80a/projects/contextual-composition/foundations/native-whole-scope-feedback/Nonblocking_Recurrent_Assembly_Reading_Route.pdf}{PDF}, \href{https://github.com/eternal-november/eos-research-programmes/blob/10e771644cfdf58c3d05e989e4b26cae7164c80a/projects/contextual-composition/foundations/native-whole-scope-feedback/Nonblocking_Recurrent_Assembly_Reading_Route.tex}{TeX}]}
locate the changed interfaces, native comparisons and stated feedback domains
mentioned in Section~\ref{sec:limits}. They are pointers to separate written
developments, not formal certification by the sequential Lean package and
not additional assumptions silently imported into this paper's theorems.

\paragraph{Authorship and proof provenance.}
GPT-6 Astra drafted this manuscript. The research and checks used multiple
AI models under human direction. This manuscript has not
undergone independent refereeing. Its prior mathematics is credited
in Appendix~\ref{app:credit} and the bibliography; the repository proofs and
formal declarations retain their distinct evidential roles. No priority
claim is inferred from the proposed terminology or from the common ambient
machinery.
\begin{thebibliography}{99}
\footnotesize\setlength{\baselineskip}{9.7pt}\setlength{\itemsep}{0pt}\setlength{\parsep}{0pt}\raggedright
\bibitem{CM} Matteo Capucci and David Jaz Myers.
\emph{Contextful Arrows via Wreath Products}.
\href{https://arxiv.org/pdf/2410.21889v2}{arXiv:2410.21889v2}.
Definitions 3.15--3.17; Appendix C, Definition C.2; Section 7; Appendix D.
\bibitem{BG} Mitchell Buckley and Richard Garner.
\emph{Composing Double Categories}.
\href{https://arxiv.org/html/2610.41551v2}{arXiv:2610.41551v2}.
Definition 6.1 and Proposition 6.3.
\bibitem{CHLL} Bastiaan Cnossen, Rune Haugseng, Tobias Lenz and Sil Linskens.
\emph{Homotopical commutative rings and bispans}.
\emph{Journal of the London Mathematical Society} 111(6), e70200, 2025.
\href{https://doi.org/10.1112/jlms.70200}{doi:10.1112/jlms.70200};
\href{https://arxiv.org/html/2403.06911v2}{arXiv:2403.06911v2}, Theorem A.
\bibitem{Lawler} Finn Lawler.
\emph{Fibrations of predicates and bicategories of relations}.
\href{https://arxiv.org/abs/1502.08017}{arXiv:1502.08017}, 2015.
\bibitem{GL} Nicola Gambino and Gabriele Lobbia.
\emph{On the formal theory of pseudomonads and pseudodistributive laws}.
\href{https://tac.mta.ca/tac/volumes/37/2/37-02abs.html}{\emph{Theory and Applications of Categories} 37(2), 14--56, 2021}.
\href{https://arxiv.org/pdf/0907.1359v3}{arXiv:0907.1359v3}.
Definition 4.1 and Appendix A.
\bibitem{Laplaza} Miquel L. Laplaza.
\href{https://www.sas.rochester.edu/mth/sites/doug-ravenel/otherpapers/laplaza.pdf}{\emph{Coherence for distributivity}}.
Lecture Notes in Mathematics 281 (1972), 29--72, Proposition 10.
\bibitem{Gluing} Contextual Composition Project.
\emph{Gluing relative corners to decorated displays}, 20 November 2026.
\mbox{[\href{https://github.com/eternal-november/eos-research-programmes/blob/a336bd6c661895d9b86827123dd38ebd8ea2a8fa/projects/contextual-composition/interaction/independent-criterion/Relative_Corners_and_Bispans_Theorem.pdf}{PDF}, \href{https://github.com/eternal-november/eos-research-programmes/blob/a336bd6c661895d9b86827123dd38ebd8ea2a8fa/projects/contextual-composition/interaction/independent-criterion/Relative_Corners_and_Bispans_Theorem.tex}{TeX}]}
\bibitem{O1} Contextual Composition Project.
\emph{General independent interaction criterion}, 1 November 2026.
\mbox{[\href{https://github.com/eternal-november/eos-research-programmes/blob/a336bd6c661895d9b86827123dd38ebd8ea2a8fa/projects/contextual-composition/interaction/independent-criterion/General_Independent_Interaction_Criterion.pdf}{PDF}, \href{https://github.com/eternal-november/eos-research-programmes/blob/a336bd6c661895d9b86827123dd38ebd8ea2a8fa/projects/contextual-composition/interaction/independent-criterion/General_Independent_Interaction_Criterion.tex}{TeX}]}
\bibitem{Comparison} Contextual Composition Project.
\emph{Which parts of an interactad are corners and bispans?} 20 November 2026.
\mbox{[\href{https://github.com/eternal-november/eos-research-programmes/blob/a336bd6c661895d9b86827123dd38ebd8ea2a8fa/projects/contextual-composition/interaction/independent-criterion/Relative_Corners_and_Bispans_Comparison.pdf}{PDF}, \href{https://github.com/eternal-november/eos-research-programmes/blob/a336bd6c661895d9b86827123dd38ebd8ea2a8fa/projects/contextual-composition/interaction/independent-criterion/Relative_Corners_and_Bispans_Comparison.tex}{TeX}]}
\bibitem{vonglehn} Tamara von Glehn. Polynomials, fibrations and distributive laws. \emph{Theory and Applications of Categories} 33(36), 1111--1144, 2018. \href{https://tac.mta.ca/tac/volumes/33/36/33-36.pdf}{Publisher PDF}.
\bibitem{gambinokock} Nicola Gambino and Joachim Kock. Polynomial functors and polynomial monads. \emph{Mathematical Proceedings of the Cambridge Philosophical Society} 154(1), 153--192, 2013. \href{https://arxiv.org/abs/0906.4931}{arXiv:0906.4931}.
\bibitem{weber} Mark Weber. Polynomials in categories with pullbacks. \emph{Theory and Applications of Categories} 30(16), 533--598, 2015. \href{https://arxiv.org/abs/1106.1983}{arXiv:1106.1983}.
\bibitem{pontos} Kate Ponto and Michael Shulman. Duality and traces for indexed monoidal categories. \emph{Theory and Applications of Categories} 26(23), 582--659, 2012. \href{https://arxiv.org/abs/1211.1555}{arXiv:1211.1555}.
\bibitem{eh} Elden Elmanto and Rune Haugseng. On distributivity in higher algebra. \emph{Compositio Mathematica} 159(11), 2321--2415, 2023. \href{https://arxiv.org/abs/2010.15722}{arXiv:2010.15722}.
\bibitem{shulman} Michael Shulman. Framed bicategories and monoidal fibrations. \emph{Theory and Applications of Categories} 20(18), 650--738, 2008. \href{https://arxiv.org/abs/0706.1286}{arXiv:0706.1286}.
\bibitem{street} Ross Street. \emph{Wreaths, mixed wreaths and twisted coactions}. Tbilisi Mathematical Journal 10(3), 1--22, 2017. \href{https://arxiv.org/abs/1610.07281}{arXiv:1610.07281}.
\bibitem{marmolejo} Francisco Marmolejo. Distributive laws for pseudomonads.
\emph{Theory and Applications of Categories} 5(5), 1999.
\href{https://tac.mta.ca/tac/volumes/1999/n5/5-05abs.html}{Publisher record}.
\bibitem{walker} Charles Walker. \emph{Presentations of pseudodistributive laws}.
\href{https://arxiv.org/abs/2102.12468v4}{arXiv:2102.12468v4}.
\bibitem{manes} Ernie Manes and Philip Mulry. Monad compositions I: General constructions and recursive distributive laws. \emph{Theory and Applications of Categories} 18(7), 172--208, 2007. \href{https://tac.mta.ca/tac/volumes/18/7/18-07abs.html}{Publisher record}.
\bibitem{fujii} Soichiro Fujii, Yun Chen Tsai, Yoàv Montacute, and Ichiro Hasuo. Monads and distributive laws in substructural contexts. \emph{LICS 2026}, LIPIcs 380, 45:1--45:28. \href{https://doi.org/10.4230/LIPIcs.LICS.2026.45}{doi:10.4230/LIPIcs.LICS.2026.45}.
\bibitem{doublefib}
Geoff Cruttwell, Michael Lambert, Dorette Pronk, and Martin Szyld.
Double fibrations. \emph{Theory and Applications of Categories} 38,
2022. \href{https://arxiv.org/abs/2205.15240}{arXiv:2205.15240}.
\bibitem{U} Contextual Composition Project.
\emph{One relative construction for context, content and effects}, 1 November 2026.
\mbox{[\href{https://github.com/eternal-november/eos-research-programmes/blob/a336bd6c661895d9b86827123dd38ebd8ea2a8fa/projects/contextual-composition/foundations/relative-construction/delivery/One_relative_construction_for_context_content_and_effects.pdf}{PDF}, \href{https://github.com/eternal-november/eos-research-programmes/blob/a336bd6c661895d9b86827123dd38ebd8ea2a8fa/projects/contextual-composition/foundations/relative-construction/delivery/One_relative_construction_for_context_content_and_effects.tex}{TeX}]}
\bibitem{Casebook} Contextual Composition Project.
\emph{Applications and Boundaries of Result-Dependent Interaction}, 1 November 2026.
\mbox{[\href{https://github.com/eternal-november/eos-research-programmes/blob/a336bd6c661895d9b86827123dd38ebd8ea2a8fa/projects/contextual-composition/applications/Applications_and_Boundaries_of_Result_Dependent_Interaction.pdf}{PDF}, \href{https://github.com/eternal-november/eos-research-programmes/blob/a336bd6c661895d9b86827123dd38ebd8ea2a8fa/projects/contextual-composition/applications/Applications_and_Boundaries_of_Result_Dependent_Interaction.tex}{TeX}]}
\bibitem{DiscreteObstruction} Contextual Composition Project.
\emph{Discrete continuation-transposition obstruction}.
\texttt{ContinuationTranspositionObstruction.no\_presentation}.
\mbox{[\href{https://github.com/eternal-november/eos-research-programmes/blob/fc2d2648058191b79b297a15161e85605cb13bb8/projects/contextual-composition/foundations/relative-construction/lean/ContextualComposition/ContinuationTranspositionObstruction.lean\#L156}{Lean},
\href{https://github.com/eternal-november/eos-research-programmes/pull/238}{verification record}]}
\bibitem{TransportIndependence} Contextual Composition Project.
\emph{Independence of Two Transport Conditions}.
\mbox{[\href{https://github.com/eternal-november/eos-research-programmes/blob/a336bd6c661895d9b86827123dd38ebd8ea2a8fa/projects/contextual-composition/interaction/independent-criterion/Transport_Conditions_Independence.pdf}{PDF}, \href{https://github.com/eternal-november/eos-research-programmes/blob/a336bd6c661895d9b86827123dd38ebd8ea2a8fa/projects/contextual-composition/interaction/independent-criterion/Transport_Conditions_Independence.tex}{TeX}]}
\bibitem{RobotTransactions} Contextual Composition Project.
\emph{Robot Transactions under Interference: Exact Compensation and the Observation Boundary}, 13 November 2026.
\mbox{[\href{https://github.com/eternal-november/eos-research-programmes/blob/a336bd6c661895d9b86827123dd38ebd8ea2a8fa/projects/contextual-composition/applications/robot-transactions/Robot_Transactions_and_Interference.pdf}{PDF}, \href{https://github.com/eternal-november/eos-research-programmes/blob/a336bd6c661895d9b86827123dd38ebd8ea2a8fa/projects/contextual-composition/applications/robot-transactions/Robot_Transactions_and_Interference.tex}{TeX}]}
\bibitem{KorthLevySilberschatz} Henry F. Korth, Eliezer Levy, and Abraham
Silberschatz. A formal approach to recovery by compensating transactions.
\emph{VLDB}, 95--106, 1990. \href{https://www.vldb.org/conf/1990/P095.PDF}{Proceedings PDF}.
\bibitem{barthe} Gilles Barthe, Thomas Espitau, Benjamin Gr\'egoire,
Justin Hsu, L\'eo Stefanesco and Pierre-Yves Strub.
\emph{Relational reasoning via probabilistic coupling}, 2015.
\href{https://arxiv.org/pdf/1509.03476v2}{arXiv:1509.03476v2}.
\bibitem{varaccawinskel} Daniele Varacca and Glynn Winskel.
\emph{Distributing probability over nondeterminism}.
\emph{Mathematical Structures in Computer Science} 16(1), 87--113, 2006.
\href{https://doi.org/10.1017/S0960129505005074}{doi:10.1017/S0960129505005074};
\href{https://www.lacl.fr/~dvaracca/papers/distrib.pdf}{author PDF}.
\end{thebibliography}
\end{document}

